% =====================================================================
% STATE NOTE (added 2026-05-02):
% This file is the formal paper draft and currently covers exp3–exp39 +
% LB ladder up to v12 (LB 0.929). Subsequent work (exp40–exp156, LB
% submissions v17–v51 culminating in v33 LB 0.932 and v51 baseline
% reset 0.917 after Kaggle test refresh) lives in:
%   - paper/exp_current.tex   running working notes (last touched
%                              2026-04-27; stops near v36)
%   - CLAUDE.md "CURRENT STATE"  authoritative ladder + invariants
%   - experiments/README.md     code-side navigation
% Author TODO before CLEF 2026-06-17 working-note deadline: port v17–v51
% ladder, pseudo-label circular-distillation finding, AudioMAE foundation
% swap (v49/v50 −0.022/−0.027), info-theoretic ceiling on optimization
% (10× RL/DPO/SFT variants ≤ BCE), and synthetic-data exploration.
% =====================================================================

\section{Experiments}
\label{sec:experiments}

This section documents our systematic experimental progression from baseline
to competitive performance on the BirdCLEF+ 2026 challenge.
All experiments use the same evaluation protocol: macro-averaged ROC-AUC on
held-out labeled soundscapes (739 segments, 75 species) as local proxy,
and Kaggle leaderboard (LB) score as the definitive metric.

% ------------------------------------------------------------------
\subsection{Experiment 3: SED Baseline}
\label{sec:exp3}

Our first competitive baseline uses Sound Event Detection (SED) with an
attention-based pooling head.

\paragraph{Setup.}
EfficientNet-B0 (NS-JFT pretrained, 5.3M params), 10-second training chunks,
mel-spectrogram input ($128 \times 224$), 2-fold ensemble.
Training: 15 epochs, AdamW (lr=$10^{-3}$), cosine schedule, batch size 64.
Augmentation: SpecAugment (freq/time masking) + spectrogram-level mixup ($\alpha=0.4$).
Dual loss: $\mathcal{L} = 0.5 \cdot \text{BCE}(\hat{y}_\text{clip}, y) + 0.5 \cdot \text{BCE}(\hat{y}_\text{frame}, y)$.

\paragraph{Inference.}
Time-shift TTA (0s and 2.5s offsets), site/hour priors from labeled soundscapes.

\paragraph{Results.}
\begin{table}[h]
\centering
\caption{Experiment 3 results (SED baseline).}
\label{tab:exp3}
\begin{tabular}{lcccc}
\hline
\textbf{Fold} & \textbf{val\_auc} & \textbf{Best Epoch} & \textbf{ss\_auc} & \textbf{LB} \\
\hline
0 & 0.9867 & 11 & 0.7252 & \multirow{2}{*}{\textbf{0.81}} \\
1 & 0.9879 & 11 & 0.7322 & \\
\hline
\end{tabular}
\end{table}

\paragraph{Analysis.}
The gap between val\_auc ($\sim$0.99) and ss\_auc ($\sim$0.73) reveals significant
domain shift between isolated training recordings and field soundscapes.
The LB score of 0.81 exceeds ss\_auc, likely due to TTA and site/hour priors
applied only at inference time.

% ------------------------------------------------------------------
\subsection{Experiment 4: Perch v2 Knowledge Distillation}
\label{sec:exp4}

\paragraph{Hypothesis.}
Google's Perch v2 bird vocalization classifier (1,536-dimensional embeddings,
$\sim$14.8K species) can provide soft supervision for 203/234 BirdCLEF species,
bridging the domain gap via teacher knowledge.

\paragraph{Setup.}
Phase 1: Extract temperature-scaled ($T=3$) soft labels from Perch v2 for all
35,549 training files (batch inference, CPU, 193.5 min).
Phase 2: Train EfficientNet-B0 student with combined loss:
$\mathcal{L} = \alpha \cdot \text{BCE}(\hat{y}, y_\text{hard}) + (1-\alpha) \cdot \text{BCE}(\hat{y}, y_\text{soft})$
where $\alpha = 0.5$. Non-bird species (31) use hard labels only.

\paragraph{Results.}
\begin{table}[h]
\centering
\caption{Experiment 4 results (Perch KD).}
\label{tab:exp4}
\begin{tabular}{lccc}
\hline
\textbf{Fold} & \textbf{val\_auc} & \textbf{ss\_auc} & \textbf{$\Delta$ vs exp3} \\
\hline
0 & 0.9502 & 0.5750 & $-0.150$ \\
1 & 0.9431 & 0.5416 & $-0.191$ \\
\hline
\end{tabular}
\end{table}

\paragraph{Analysis.}
KD from Perch \textbf{degraded} soundscape performance substantially.
val\_auc was lower (0.95 vs 0.99), and ss\_auc dropped from 0.73 to 0.55--0.58.
The failure likely stems from domain mismatch: Perch was trained on global bird
recordings, not Pantanal field soundscapes with overlapping multi-species audio.
The soft labels pulled the student toward Perch's species priors rather than
learning the target distribution.

\textbf{Finding:} Naive KD from a general-purpose audio classifier hurts
performance when the teacher's training domain differs significantly from
the target evaluation domain.

% ------------------------------------------------------------------
\subsection{Experiment 5: Self-Training v1}
\label{sec:exp5}

\paragraph{Hypothesis.}
Pseudo-labeling 10,592 unlabeled soundscapes with the exp3 teacher and
retraining should improve soundscape-domain performance.

\paragraph{Setup.}
Teacher: exp3 2-fold ensemble. Pseudo-label threshold: 0.5 (confidence).
Power transform ($p=0.7$) on pseudo-label probabilities.
Student initialized from exp3 weights, fine-tuned with
lr=$5 \times 10^{-4}$, 15 epochs, pseudo sample weight 0.5.

\paragraph{Results.}
Only 935/10,592 (8.8\%) soundscapes exceeded the confidence threshold,
yielding 5,610 pseudo training samples.

\begin{table}[h]
\centering
\caption{Experiment 5 results (self-training v1).}
\label{tab:exp5}
\begin{tabular}{lcccc}
\hline
\textbf{Fold} & \textbf{val\_auc} & \textbf{Best Epoch} & \textbf{ss\_auc} & \textbf{$\Delta$ vs exp3} \\
\hline
0 & 0.9963 & 1 & 0.7161 & $-0.009$ \\
1 & 0.9949 & 1 & 0.7482 & $+0.016$ \\
\hline
\end{tabular}
\end{table}

\paragraph{Analysis.}
The best checkpoint for both folds was epoch 1 (immediately after initialization
from exp3), indicating that continued fine-tuning degraded performance.
Key issues: (1) learning rate too high for fine-tuning ($5 \times 10^{-4}$),
(2) only 8.8\% of soundscapes produced pseudo-labels (threshold too strict),
(3) 15 epochs caused overfitting from the very start.

Fold 1 showed marginal improvement (+0.016 ss\_auc), suggesting pseudo-labels
have potential but require careful hyperparameter tuning.

\textbf{Finding:} Self-training from a pre-trained checkpoint requires very low
learning rate and minimal epochs. The pseudo-label threshold of 0.5 was too
conservative, capturing only 8.8\% of available unlabeled data.

% ------------------------------------------------------------------
\subsection{Experiment 6: Backbone Upgrade + Extended Context}
\label{sec:exp6}

\paragraph{Hypothesis.}
Scaling the backbone from EfficientNet-B0 (5.3M) to B3 (12M) increases model
capacity, and extending training chunks from 10s to 20s provides richer
temporal context for species identification.

\paragraph{Setup.}
Backbone: \texttt{tf\_efficientnet\_b3.ns\_jft\_in1k} (12M params).
Training chunks: 20 seconds. Image size: $128 \times 300$.
Batch size reduced to 32 (GPU memory). Otherwise identical to exp3.

\paragraph{Results.}

\begin{table}[h]
\centering
\caption{Experiment 6 results: backbone and duration ablation.}
\label{tab:exp6}
\begin{tabular}{llcccc}
\hline
\textbf{Variant} & \textbf{Backbone} & \textbf{Duration} & \textbf{Fold} & \textbf{val\_auc} & \textbf{ss\_auc} \\
\hline
exp6 (20s)  & B3 (13.8M) & 20s & 0 & 0.9840 & 0.6344 \\
exp6b (10s) & B3 (13.8M) & 10s & 0 & 0.9845 & 0.7206 \\
exp6b (10s) & B3 (13.8M) & 10s & 1 & 0.9869 & 0.6949 \\
\hline
exp3 (ref)  & B0 (5.3M)  & 10s & 0 & 0.9867 & 0.7252 \\
exp3 (ref)  & B0 (5.3M)  & 10s & 1 & 0.9879 & 0.7322 \\
\hline
\end{tabular}
\end{table}

\paragraph{Analysis.}
Two key findings emerge:

\textbf{(1) Training duration mismatch hurts evaluation.}
The 20s variant achieved comparable val\_auc (0.984) but significantly lower
ss\_auc (0.634 vs 0.725). Since evaluation uses 5-second segments, the
spectrogram resize from 20s to a fixed image size creates very different
temporal patterns than seen during training. This 4:1 temporal compression
ratio degrades learned features at inference time.

\textbf{(2) Backbone capacity is not the bottleneck.}
EfficientNet-B3 (13.8M params, 2.6$\times$ B0) produced nearly identical
ss\_auc to B0 when using the same 10s duration. The domain shift between
isolated training recordings and multi-species field soundscapes---not model
capacity---is the primary performance limiter. Future experiments should focus
on closing this domain gap rather than scaling model size.

% ------------------------------------------------------------------
\subsection{Experiment 7: Self-Training v2 — Domain Adaptation}
\label{sec:exp7}

\paragraph{Hypothesis.}
The domain shift between \texttt{train\_audio} (isolated recordings) and
soundscapes (multi-species field recordings) is the primary performance limiter.
Directly injecting soundscape data into training and expanding pseudo-labels
should close this gap.

\paragraph{Key changes from exp5 (self-training v1).}
\begin{enumerate}
    \item Labeled soundscapes included in training (50/50 train/eval split by file)
    \item Pseudo-label threshold lowered: $0.5 \to 0.3$ (captures 28.2\% vs 8.8\%)
    \item Learning rate: $5 \times 10^{-4} \to 10^{-5}$ (fine-tuning, not retraining)
    \item Epochs: $15 \to 5$ (prevents overfitting)
    \item Pseudo sample weight: $0.5 \to 0.3$
    \item Best checkpoint selected by ss\_auc (not val\_auc)
\end{enumerate}

\paragraph{Results.}
Training data: 54,025 audio + 359 labeled soundscape + 17,892 pseudo = 72,276 total.

\begin{table}[h]
\centering
\caption{Experiment 7 results (self-training v2).}
\label{tab:exp7}
\begin{tabular}{lcccc}
\hline
\textbf{Fold} & \textbf{ss\_auc (held-out)} & \textbf{full\_ss} & \textbf{Best Epoch} & \textbf{$\Delta$ vs exp3} \\
\hline
0 & 0.8130 & 0.7810 & 4 & $+0.056$ \\
1 & 0.7902 & 0.8054 & 4 & $+0.073$ \\
\hline
\end{tabular}
\end{table}

\paragraph{Analysis.}
This is the largest single improvement in our experimental series.
Three factors contributed:

\textbf{(1) Soundscape training data.}
Including even 359 labeled soundscape segments (0.5\% of training data)
provided direct exposure to the target domain's acoustic characteristics.

\textbf{(2) Lower pseudo-label threshold.}
Reducing from 0.5 to 0.3 tripled the pseudo-labeled data (17,892 vs 5,610),
providing substantially more domain-relevant training signal.

\textbf{(3) Conservative fine-tuning.}
The combination of very low learning rate ($10^{-5}$) and few epochs (5)
preserved the base model's learned features while adapting to the new data.
Best checkpoints were at epoch 4, confirming the narrow useful training window.

\textbf{Finding:} Domain adaptation through in-domain data inclusion and
conservative pseudo-label fine-tuning is far more effective than model scaling
(exp6) or knowledge distillation from out-of-domain teachers (exp4).

% ------------------------------------------------------------------
\subsection{Experiment 8: Focal Loss + Label Smoothing}
\label{sec:exp8}

\paragraph{Hypothesis.}
Focal loss ($\gamma=2$) down-weights easy negatives and focuses gradient on
hard examples (rare species, ambiguous soundscapes), while label smoothing
($\epsilon=0.05$) improves probability calibration for AUC evaluation.

\paragraph{Setup.}
Same pipeline as exp7: EfficientNet-B0, initialized from exp3, combined
training data (train\_audio + labeled soundscapes + pseudo-labels from exp3
teacher at threshold 0.3), LR=$10^{-5}$, 5 epochs, cosine schedule.

\textbf{Change from exp7:} BCE $\to$ focal loss with
$\mathcal{L}_\text{focal} = (1-p_t)^\gamma \cdot \text{BCE}(\hat{y}, y')$
where $y' = y(1-\epsilon) + 0.5\epsilon$ (label smoothing).
Applied to both clipwise and framewise heads.

\paragraph{Results.}

\begin{table}[h]
\centering
\caption{Experiment 8 results (focal loss + label smoothing).}
\label{tab:exp8}
\begin{tabular}{lccccc}
\hline
\textbf{Fold} & \textbf{ss\_auc} & \textbf{full\_ss} & \textbf{Best Epoch} & \textbf{$\Delta$ vs exp7} \\
\hline
0 & 0.8282 & 0.8149 & 5 & $+0.034$ \\
1 & 0.8051 & 0.8242 & 5 & $+0.019$ \\
\hline
\end{tabular}
\end{table}

\paragraph{Analysis.}
Focal loss consistently improved over exp7's BCE baseline across both folds,
with full\_ss gains of +0.034 (fold 0) and +0.019 (fold 1).
Three contributing mechanisms:

\textbf{(1) Hard example mining via focal loss.}
With $\gamma=2$, easy negatives (the vast majority of the 234 classes per sample)
receive near-zero gradient. This effectively concentrates learning on the
$\sim$5--10 confusable species per sample, particularly benefiting rare taxa.

\textbf{(2) Label smoothing for AUC calibration.}
Smoothing targets from $\{0,1\}$ to $\{0.025, 0.975\}$ prevents the model from
producing overconfident predictions. Since AUC depends on ranking rather than
absolute calibration, the reduced overconfidence improves separation between
true positives and hard negatives.

\textbf{(3) Extended training window.}
Both folds achieved their best checkpoint at epoch 5 (vs epoch 4 in exp7),
suggesting that focal loss's implicit curriculum (easy$\to$hard) enables
productive training for longer without overfitting.

\textbf{Finding:} Focal loss with label smoothing provides a complementary
improvement orthogonal to domain adaptation (exp7). The combination of
self-training + focal loss yields full\_ss 0.81--0.82, a cumulative
improvement of $+0.09$ over the exp3 baseline.

% ------------------------------------------------------------------
\subsection{Experiment 9: Enhanced Augmentation + 2nd Round Self-Training}
\label{sec:exp9}

\paragraph{Hypothesis.}
Iterative self-training with an improved teacher (exp8) produces higher-quality
pseudo-labels, while stronger augmentation (aggressive SpecAugment, random gain,
Gaussian noise) acts as noise injection that improves robustness.

\paragraph{Key changes from exp8.}
\begin{enumerate}
    \item Teacher upgraded: exp3 $\to$ exp8 ensemble for pseudo-labeling
    \item Pseudo threshold lowered: $0.3 \to 0.2$ (exp8 teacher is more confident)
    \item Init from exp8 weights (not exp3)
    \item Learning rate: $10^{-5} \to 5 \times 10^{-6}$ (2nd round fine-tuning)
    \item Mixup $\alpha$: $0.4 \to 0.5$
    \item SpecAugment: up to 3 freq masks ($w \leq 30$), 3 time masks ($w \leq 40$)
    \item Random gain ($\pm 6$ dB) and Gaussian noise injection ($p=0.3$)
\end{enumerate}

\paragraph{Results.}
Teacher upgrade yielded 97.4\% pseudo-label pass rate (vs 79.1\% with exp3
teacher at threshold 0.3), producing 17,430 pseudo samples.
Combined training set: 71,814 samples.

\begin{table}[h]
\centering
\caption{Experiment 9 results (enhanced aug + 2nd round ST).}
\label{tab:exp9}
\begin{tabular}{lccccc}
\hline
\textbf{Fold} & \textbf{ss\_auc} & \textbf{full\_ss} & \textbf{Best Epoch} & \textbf{$\Delta$ vs exp8} \\
\hline
0 & 0.8403 & 0.8262 & 5 & $+0.011$ \\
1 & 0.8124 & 0.8301 & 3 & $+0.006$ \\
\hline
\end{tabular}
\end{table}

\paragraph{Analysis.}
Two orthogonal improvements contributed:

\textbf{(1) Iterative self-training with better teacher.}
The exp8 teacher's pseudo-label pass rate jumped from 79.1\% to 97.4\%,
adding $\sim$3K more pseudo-labeled samples. More importantly, the
pseudo-labels themselves are higher quality, as the exp8 teacher already
benefits from focal loss's hard-example emphasis.

\textbf{(2) Stronger augmentation as regularization.}
Enhanced SpecAugment (3$\times$ masks, wider), random gain, and Gaussian noise
serve as implicit noise injection for self-training. This follows the
Noisy Student paradigm~\cite{xie2020self}: adding noise during student
training prevents the student from merely memorizing teacher predictions,
forcing it to learn more robust features.

\textbf{(3) Diminishing returns observed.}
The improvement from exp8$\to$exp9 ($+0.006$--$0.011$) is smaller than
exp7$\to$exp8 ($+0.019$--$0.034$), suggesting that the easy gains from
domain adaptation and loss function optimization have been captured.
Further improvement likely requires architectural diversity or
post-processing rather than continued iteration on the same recipe.

\textbf{Finding:} Iterative self-training with progressive teacher upgrades
yields consistent but diminishing gains. The full pipeline (domain adaptation +
focal loss + 2nd round ST + strong augmentation) achieves full\_ss
0.83, a cumulative $+0.10$ over the exp3 baseline.

% ------------------------------------------------------------------
\subsection{Experiment 10: Ensemble Optimization + Submission}
\label{sec:exp10}

\paragraph{Hypothesis.}
Combining models from diverse training recipes (exp7: self-training v2,
exp8: focal loss, exp9: enhanced augmentation) creates an ensemble that
captures complementary patterns, improving over any single experiment.

\paragraph{Setup.}
Inference-only experiment: load all 8 model checkpoints (exp3/7/8/9 $\times$ 2
folds) and evaluate ensemble combinations on the full soundscape evaluation set
(739 segments, 75 species). Tested: individual models, within-experiment
averaging, cross-experiment combinations, weighted ensembles, and power transform
post-processing.

\paragraph{Results.}

\begin{table}[h]
\centering
\caption{Experiment 10: ensemble comparison on full\_ss.}
\label{tab:exp10}
\begin{tabular}{lcc}
\hline
\textbf{Configuration} & \textbf{full\_ss} & \textbf{$\Delta$ vs best single} \\
\hline
Best single (exp9 f1)        & 0.8302 & --- \\
exp9 ensemble (2 folds)      & 0.8450 & $+0.015$ \\
exp8+exp9 (4 models)         & 0.8452 & $+0.015$ \\
\textbf{exp7+exp8+exp9 (6)}  & \textbf{0.8466} & $+\mathbf{0.016}$ \\
All models (8, incl exp3)    & 0.8463 & $+0.016$ \\
Weighted (3:2:1:0.5)         & 0.8464 & $+0.016$ \\
\hline
\end{tabular}
\end{table}

\paragraph{Analysis.}
Key findings from ensemble optimization:

\textbf{(1) Fold averaging provides the largest gain.}
Going from a single fold to 2-fold average (exp9 alone: 0.830 $\to$ 0.845)
accounts for most of the ensemble benefit ($+0.015$). Cross-experiment
diversity adds only $+0.001$--$0.002$ on top.

\textbf{(2) Weak models hurt the ensemble.}
Including exp3 (full\_ss $\sim$0.73) slightly degraded the 6-model ensemble
from 0.8466 to 0.8463. The optimal strategy is to include only models within
$\sim$0.05 of the best.

\textbf{(3) Post-processing yielded no improvement.}
Power transform (tested at $p \in \{0.5, \ldots, 0.9\}$) on non-bird species
had zero effect, as the evaluation set contains few non-bird true positives.
Weighted ensembles marginally underperformed uniform averaging.

\textbf{(4) Adaptive inference for CPU budget.}
With $\sim$600 test files, 6 models, and 3 TTA shifts, inference would require
$\sim$160 minutes on CPU, exceeding the 90-minute budget. The submission
notebook dynamically selects model count and TTA based on test set size.

\paragraph{Leaderboard ablation.}
Submitting six inference variants to the public leaderboard revealed a
\textbf{critical disconnect} between local proxy and LB score.

\begin{table}[h]
\centering
\caption{Leaderboard ablation: local full\_ss vs public LB.}
\label{tab:lb-ablation}
\begin{tabular}{llcc}
\hline
\textbf{Variant} & \textbf{Configuration} & \textbf{Local full\_ss} & \textbf{Public LB} \\
\hline
exp3 baseline     & B0 2-fold, TTA, priors            & 0.732 & \textbf{0.811} \\
exp10 v2          & 6-model (exp7+8+9), TTA, priors   & 0.847 & 0.797 \\
v3                & exp9-only, TTA, priors             & 0.845 & 0.797 \\
v4                & 6-model, TTA, \textbf{no priors}   & 0.847 & 0.790 \\
v5                & exp9-only, no TTA, no priors        & 0.830 & 0.785 \\
v6                & 6-model, TTA, priors, smoothing     & ---   & 0.803 \\
\hline
\end{tabular}
\end{table}

\paragraph{Critical finding: local ss\_auc anti-correlates with LB.}
Models that improved local full\_ss from 0.73 to 0.85 \textbf{degraded} LB from
0.811 to 0.785--0.803. This reveals fundamental flaws in our experimental setup:

\textbf{(1) Self-training overfit to train\_soundscapes distribution.}
Experiments 7--9 used pseudo-labels from train\_soundscapes, which biased the
model toward the acoustic characteristics of those specific recordings (sites,
times, species mixtures). The test soundscapes follow a different distribution
from the labeled training soundscapes, so this ``domain adaptation'' actually
\emph{narrowed} the model's generalization instead of broadening it.

\textbf{(2) Pseudo-label noise amplification across rounds.}
Each self-training round (exp3 $\to$ exp8 $\to$ exp9) reinforced the teacher's
errors. Systematic biases (e.g., species confused by exp3) became entrenched
rather than corrected, compounding over iterations.

\textbf{(3) The exp3 baseline generalizes better precisely because it is naive.}
Trained purely on isolated train\_audio recordings, exp3 learns species-specific
features without being biased toward any particular soundscape distribution.
This generality transfers better to unseen test soundscapes.

\textbf{(4) Inference-time techniques provide genuine gains.}
In contrast to training-time changes, the LB ablation confirms:
TTA contributes $+0.012$ (v5 $\to$ v3), site/hour priors $+0.007$--$0.013$
(v4 $\to$ v6), and temporal smoothing $+0.006$ (v2 $\to$ v6).
These are additive and do not overfit.

\textbf{Lesson:} When labeled evaluation data is not representative of the
test distribution, optimizing a local proxy can be counterproductive.
Future experiments must validate against LB directly, and training-time
modifications (self-training, loss changes) must be treated as hypotheses
until confirmed on LB.

% ------------------------------------------------------------------
\subsection{Experiment 11: 5-Second Chunk Alignment}
\label{sec:exp11}

\paragraph{Hypothesis.}
The 10-second training chunks in exp3 create a 2:1 temporal resolution mismatch
with the 5-second evaluation windows: mel-spectrograms are resized to a fixed
image size during both training and inference, but a 10s spectrogram compressed
into the same spatial dimensions as a 5s spectrogram produces very different
temporal patterns. Aligning training to 5s chunks should eliminate this distortion.

\paragraph{Key changes from exp3.}
\begin{enumerate}
    \item Training chunk duration: $10\text{s} \to 5\text{s}$
    \item Max chunks per file: $2 \to 4$ (maintains similar coverage of each recording)
    \item Training samples: 104,134 (vs 54,025 with 10s chunks)
    \item Precomputed mel-spectrogram cache to minimize CPU load
\end{enumerate}
All other settings (backbone, optimizer, augmentation, loss) remain identical to exp3.

\paragraph{Results.}
\begin{table}[h]
\centering
\caption{Experiment 11 results (5s chunk alignment).}
\label{tab:exp11}
\begin{tabular}{lcccc}
\hline
\textbf{Fold} & \textbf{val\_auc} & \textbf{Best Epoch} & \textbf{ss\_auc} & \textbf{$\Delta$ vs exp3} \\
\hline
0 & 0.9862 & 9  & 0.7715 & $+0.046$ \\
1 & 0.9857 & 10 & 0.7503 & $+0.018$ \\
\hline
Average & 0.9860 & --- & 0.7609 & $+0.032$ \\
\hline
\end{tabular}
\end{table}

\paragraph{Analysis.}
The 5s alignment provides a consistent and substantial improvement (+0.032
average ss\_auc) with no architectural changes. This confirms that temporal
resolution mismatch was a significant performance bottleneck: all competitive
public notebooks use 5s training chunks matching the evaluation window.
The nearly doubled training set (104K vs 54K samples) from finer chunking
may also contribute, as the model sees more diverse 5s views of each recording.

\textbf{Finding:} Aligning training chunk duration to the evaluation window
is a prerequisite for competitive performance. The 10s$\to$5s change alone
accounts for a +0.032 ss\_auc gain with zero additional complexity.

% ------------------------------------------------------------------
\subsection{Experiment 12: Inference Optimization}
\label{sec:exp12}

\paragraph{Hypothesis.}
Inference-time techniques---overlapping TTA, temperature scaling, and temporal
smoothing---can improve prediction quality without retraining. We conduct a
grid search over these axes using exp3 and exp11 weights.

\paragraph{Setup.}
Grid search over: (1) TTA mode: standard (non-overlapping 5s) vs overlap
(2.5s stride, 23 windows per 60s file); (2) temperature scaling:
$T \in \{1.0, 1.15, 1.3, 1.5\}$; (3) Gaussian temporal smoothing:
$\sigma \in \{0, 0.5, 1.0, 1.5\}$. Evaluated on the full soundscape set
(739 segments, 75 species).

\paragraph{Results.}
\begin{table}[h]
\centering
\caption{Experiment 12: top inference configurations (ss\_auc on full soundscape set).}
\label{tab:exp12}
\begin{tabular}{llccc}
\hline
\textbf{Model} & \textbf{TTA} & \textbf{$T$} & \textbf{$\sigma$} & \textbf{ss\_auc} \\
\hline
exp11 & overlap & 1.0 & 1.0 & \textbf{0.8015} \\
exp11 & overlap & 1.0 & 0.5 & 0.8004 \\
exp11 & overlap & 1.15 & 1.0 & 0.8000 \\
exp11 & standard & 1.0 & 1.0 & 0.7951 \\
exp3+exp11 & overlap & 1.0 & 1.0 & 0.7984 \\
exp3 & overlap & 1.0 & 1.0 & 0.7792 \\
exp3 & standard & 1.0 & 0.0 & 0.7502 \\
\hline
\end{tabular}
\end{table}

\paragraph{Analysis.}
Three clear findings emerge from the 96-configuration grid search:

\textbf{(1) Temperature scaling has no effect.}
At every fixed (model, TTA, $\sigma$) combination, varying $T$ from 1.0 to 1.5
changes ss\_auc by $<$0.002. Since our predictions are already sigmoids
(not softmax), temperature scaling merely compresses probabilities toward 0.5,
which does not affect AUC ranking. $T=1.0$ is universally optimal.

\textbf{(2) Gaussian temporal smoothing consistently helps.}
$\sigma=1.0$ provides $+0.009$ average improvement over no smoothing.
This confirms that adjacent 5s windows share species presence, and averaging
reduces segment-boundary noise.

\textbf{(3) Overlap TTA provides +0.015.}
The 2.5s-stride sliding window (23 predictions per file vs 12) captures
species vocalizations that straddle segment boundaries. Combined with
smoothing, this yields a +0.023 total inference-time gain.

\textbf{(4) Cross-architecture ensemble hurts.}
Adding exp3 (10s-trained) to exp11 (5s-trained) degrades from 0.8015 to 0.7984.
The 10s model's distorted temporal features add noise rather than diversity.

\textbf{Finding:} Optimal inference config is overlap TTA + $\sigma$=1.0 +
$T$=1.0, yielding 0.8015 from exp11 weights alone---a +0.051 gain over
exp3's raw inference (0.7502).

% ------------------------------------------------------------------
\subsection{Experiment 13: Multi-Backbone Exploration}
\label{sec:exp13}

\paragraph{Hypothesis.}
Backbone diversity is critical for ensemble robustness (per 2025 winning
solutions). We evaluate three alternative architectures---EfficientNet-B2,
ConvNeXt-V2 Nano, and NFNet-L0---using the 5s training recipe from exp11.

\paragraph{Results.}
\begin{table}[h]
\centering
\caption{Experiment 13 results (backbone comparison, 5s chunks).}
\label{tab:exp13}
\begin{tabular}{lcccc}
\hline
\textbf{Backbone} & \textbf{Params} & \textbf{Fold} & \textbf{ss\_auc} & \textbf{Status} \\
\hline
EfficientNet-B0 (exp11) & 5.3M & 0/1 & 0.771/0.750 & Baseline \\
EfficientNet-B2 & 9.1M & 0 & 0.7090 & $-0.063$ \\
EfficientNet-B2 & 9.1M & 1 & 0.7108 & $-0.039$ \\
ConvNeXt-V2 Nano & 15.6M & 0 & --- & Failed \\
NFNet-L0 & 35.1M & 0 & --- & Failed \\
\hline
\end{tabular}
\end{table}

\paragraph{Analysis.}
All three alternatives performed worse than the B0 baseline:

\textbf{(1) EfficientNet-B2} completed training but underperformed B0 by
0.04--0.06 ss\_auc despite 1.7$\times$ more parameters. With only 5s of audio
($128 \times 156$ mel-spectrogram), the input resolution is too low to benefit
from deeper features. The additional capacity leads to overfitting on
train\_audio without improving soundscape generalization.

\textbf{(2) ConvNeXt-V2 Nano} failed entirely: val\_auc remained at 0.50
(random chance) throughout training. ConvNeXt-V2 uses LayerNorm which
is incompatible with single-channel mel-spectrogram input (\texttt{in\_chans=1})
in our SED pipeline without architecture-specific normalization adaptation.

\textbf{(3) NFNet-L0} produced NaN losses immediately. NFNet's Scaled Weight
Standardization is incompatible with mixed-precision training (AMP) at
lr=$10^{-3}$, causing numerical overflow.

\textbf{Finding:} EfficientNet-B0 remains optimal for this task and pipeline.
Backbone diversity requires architecture-specific adaptation (normalization,
learning rate, AMP settings) rather than drop-in replacement. Future ensemble
diversity should come from training recipe variation rather than backbone changes.

% ------------------------------------------------------------------
\subsection{Experiment 14: Mel-Level Data Augmentation}
\label{sec:exp14}

\paragraph{Hypothesis.}
Stronger data augmentation at the mel-spectrogram level---background noise
mixing, random gain, and Gaussian noise---can improve generalization to
field soundscapes without CPU-intensive waveform processing during training.

\paragraph{Setup.}
All augmentations operate on precomputed mel-spectrograms (reusing exp11's
cache), eliminating librosa overhead during training:
\begin{enumerate}
    \item \textbf{Background mixing} ($p=0.5$): mix with precomputed mel-spectrograms
    from 2025 train\_soundscapes (1,200 mels from 300 files, computed once).
    Mix ratio sampled from Beta(5,2), mean $\approx$0.7 signal / 0.3 background.
    \item \textbf{Random gain} ($p=0.5$): multiply mel by uniform(0.8, 1.2)
    \item \textbf{Gaussian noise} ($p=0.3$): additive $\mathcal{N}(0, 0.02)$
    \item \textbf{Enhanced SpecAugment}: up to 3 frequency masks ($w \leq 30$),
    3 time masks ($w \leq 30$)
\end{enumerate}

\paragraph{Results.}
\begin{table}[h]
\centering
\caption{Experiment 14 results (mel-level augmentation).}
\label{tab:exp14}
\begin{tabular}{lcccc}
\hline
\textbf{Fold} & \textbf{val\_auc} & \textbf{Best Epoch} & \textbf{ss\_auc} & \textbf{$\Delta$ vs exp11} \\
\hline
0 & 0.9871 & 13 & 0.7820 & $+0.011$ \\
1 & 0.9856 & 13 & 0.7681 & $+0.018$ \\
\hline
Average & 0.9864 & --- & 0.7751 & $+0.014$ \\
\hline
\end{tabular}
\end{table}

\paragraph{Analysis.}
The augmentation provides a consistent +0.014 average improvement, with two
noteworthy effects:

\textbf{(1) Extended training window.}
Best checkpoints shifted from epochs 9--10 (exp11) to epoch 13 (exp14).
In exp11, ss\_auc peaked early and declined due to overfitting on clean
train\_audio. The added augmentation noise acts as regularization, allowing
productive training for 3--4 additional epochs before overfitting onset.

\textbf{(2) Background mixing bridges domain gap.}
Mixing training audio with real soundscape backgrounds at the mel level
simulates the multi-source acoustic environment of field recordings. This
is computationally efficient: 1,200 background mels are precomputed once
(one-time cost), and mixing is a simple element-wise operation during training.

\textbf{(3) CPU efficiency.}
By operating entirely on precomputed mels, training CPU usage is minimal---no
librosa decoding, no waveform-level processing. All augmentation is GPU-friendly
tensor operations.

\textbf{Finding:} Mel-level augmentation with background mixing provides a
meaningful +0.014 ss\_auc gain while extending the useful training window.
The approach is orthogonal to architecture and inference optimizations.

% ------------------------------------------------------------------
\subsection{Experiment 15: Self-Training with 2025 Winner Recipe}
\label{sec:exp15}

\paragraph{Hypothesis.}
The 2025 BirdCLEF winner's self-training recipe---which was the single most
impactful technique that year---should improve our 5s baseline when properly
implemented with noise injection (always-on mixup, stochastic depth) and
power-transformed pseudo-labels.

\paragraph{Setup.}
Faithful reproduction of the 2025 1st-place self-training recipe:
\begin{enumerate}
    \item Teacher: exp11 2-fold ensemble
    \item Pseudo-labels: power transform ($p=0.7$) on teacher predictions.
    62,040 pseudo training samples (48.8\% pass rate)
    \item Always-on mixup: every training sample mixed with a weighted-sampled
    pseudo sample (\texttt{mixup\_ratio}=1.0)
    \item Stochastic depth: \texttt{drop\_path\_rate}=0.15
    \item Sampling weight proportional to sum of pseudo-label probabilities
    \item Training: 5 epochs, lr=$10^{-3}$, from scratch (not fine-tuning)
    \item Checkpoint selection by ss\_auc (not val\_auc)
\end{enumerate}

\paragraph{Results.}
\begin{table}[h]
\centering
\caption{Experiment 15 results (2025 winner self-training recipe).}
\label{tab:exp15}
\begin{tabular}{lccccc}
\hline
\textbf{Fold} & \textbf{Best Epoch} & \textbf{val\_auc} & \textbf{ss\_auc} & \textbf{$\Delta$ vs exp11} \\
\hline
0 & 1 & 0.9823 & 0.7569 & $-0.015$ \\
1 & 2 & 0.9809 & 0.6908 & $-0.060$ \\
\hline
\end{tabular}
\end{table}

\paragraph{Epoch-level analysis.}
\begin{table}[h]
\centering
\caption{Experiment 15 ss\_auc trajectory (fold 0).}
\label{tab:exp15-trajectory}
\begin{tabular}{lccccc}
\hline
\textbf{Epoch} & 1 & 2 & 3 & 4 & 5 \\
\hline
ss\_auc (fold 0) & \textbf{0.757} & 0.750 & 0.746 & 0.744 & 0.745 \\
ss\_auc (fold 1) & 0.691 & \textbf{0.691} & 0.685 & 0.679 & 0.683 \\
\hline
\end{tabular}
\end{table}

\paragraph{Analysis.}
Self-training \textbf{degraded} performance relative to the exp11 baseline,
with ss\_auc peaking at epoch 1--2 and declining thereafter. This is our third
failed attempt at self-training (after exp5 and exp7--9), despite following the
2025 winner's exact recipe. Three factors explain the persistent failure:

\textbf{(1) Missing LB validation loop.}
The 2025 winners used leaderboard scores to select which self-training
iterations to keep. Without this external validation signal, we cannot
distinguish between pseudo-labels that genuinely improve generalization and
those that overfit to train\_soundscapes. Our ss\_auc proxy is unreliable
for this purpose (as demonstrated by exp7--9's LB anti-correlation).

\textbf{(2) Teacher quality ceiling.}
Our exp11 teacher (ss\_auc=0.76) is significantly weaker than the 2025
winners' teachers (which had LB $>$0.85). Weak teachers produce noisy
pseudo-labels that corrupt the student more than they help, especially
on a diverse 234-species task.

\textbf{(3) Distribution shift from pseudo-labels.}
The pseudo-labeled train\_soundscapes represent a specific acoustic
distribution (particular sites, times, species mixtures). Training on this
data shifts the model toward that distribution, which may not match the
hidden test soundscapes---the same failure mode observed in exp7--9.

\textbf{Finding:} Self-training is not viable without (a) a strong teacher
model ($>$0.85 LB) and (b) leaderboard validation per iteration. Both
prerequisites require significant compute investment and careful orchestration.
For our current pipeline, augmentation-based approaches (exp14) provide more
reliable gains.

% ------------------------------------------------------------------
\subsection{Summary: Experiments 11--15}
\label{sec:exp11-15-summary}

Table~\ref{tab:exp11-15-summary} summarizes the cumulative effect of
experiments 11--15, starting from the exp3 baseline.

\begin{table}[h]
\centering
\caption{Summary of experiments 11--15 (average ss\_auc across folds).}
\label{tab:exp11-15-summary}
\begin{tabular}{llcc}
\hline
\textbf{Exp} & \textbf{Key Change} & \textbf{Avg ss\_auc} & \textbf{$\Delta$ vs exp3} \\
\hline
exp3  & SED baseline (10s chunks) & 0.729 & --- \\
exp11 & 5s chunk alignment & 0.761 & $+0.032$ \\
exp12 & + overlap TTA + $\sigma$=1.0 (inference) & 0.802 & $+0.073$ \\
exp13 & Multi-backbone (B2/ConvNeXt/NFNet) & 0.710 & $-0.019$ \\
exp14 & + mel-level augmentation & 0.775 & $+0.046$ \\
exp15 & Self-training (2025 recipe) & 0.724 & $-0.005$ \\
\hline
\end{tabular}
\end{table}

The most impactful changes are training-eval alignment (exp11: $+0.032$) and
inference optimization (exp12: $+0.041$ on top of exp11). Together, these
bring the local proxy from 0.729 to 0.802 without modifying the model
architecture or training recipe beyond chunk duration. Mel-level augmentation
(exp14) provides an additional $+0.014$ that should compose with inference
optimizations.

The two negative results---backbone scaling (exp13) and self-training
(exp15)---confirm findings from our earlier experiments: the performance
bottleneck is not model capacity, and self-training without LB validation
is unreliable. Future experiments should focus on (a) submitting exp14 +
exp12's inference config to validate LB improvement, and (b) exploring
training-time techniques that do not depend on pseudo-labels from
train\_soundscapes.

% ------------------------------------------------------------------
\subsection{Competitive Landscape Analysis}
\label{sec:landscape}

To inform future experiments, we surveyed the public leaderboard and
9 public Kaggle notebooks as of April 17, 2026.

\paragraph{Leaderboard snapshot.}
Top-1 scores 0.952, top-10 all $\geq 0.944$, top-20 $\geq 0.938$.
Our best submission (0.811) would rank approximately in the bottom third.
The gap to a medal-competitive score ($\sim$0.93) is $+0.12$.

\paragraph{Public notebook survey.}

\begin{table}[h]
\centering
\caption{Summary of public BirdCLEF+ 2026 inference notebooks.}
\label{tab:landscape}
\small
\begin{tabular}{p{3.2cm}p{2.8cm}p{2.2cm}p{4.5cm}}
\hline
\textbf{Notebook} & \textbf{Backbone} & \textbf{Chunk / Image} & \textbf{Key techniques} \\
\hline
ConvNeXt-V2 + GEM & ConvNeXt-V2 tiny & 5s / 128$\times$W & Overlap TTA (2.5s stride), temporal smoothing, file-level max prior \\
EfficientNetV2 & EfficientNetV2-B0 & 5s / 256$\times$256 & 256 mel bins, ThreadPool parallel inference \\
BirdTransform & CNN + Transformer (RoPE, SwiGLU) & 7s / 192 mels & PERCH embedding fusion ($\mu$=0.5), per-class threshold optimization \\
VK EfficientNet & EfficientNet-B0 & 5s / 256$\times$256 & Multi-model ensemble, flip TTA, fmin=50 fmax=14k \\
Perch v2 & Google Perch v2 & overlap / 5s & Foundation model (15k species), Gaussian temporal smoothing ($\sigma$=1.0) \\
Kailatha B2 + SED & EfficientNet-B2 & 5s / 128$\times$312 & SED attention pool, temperature scaling ($T$=1.5) \\
MonkeyClaw & Custom 3-layer CNN & 5s / 128$\times$128 & GroupKFold by site, temperature scaling ($T$=1.15), prior fusion \\
Kahofu & None (statistical) & --- & Rating-weighted species frequency prior \\
\hline
\end{tabular}
\end{table}

\paragraph{Key observations.}
\begin{enumerate}
    \item \textbf{All competitive solutions use 5s chunks}, matching the
    evaluation window. Our 10s training chunks require resize at inference,
    introducing temporal distortion. Aligning train/eval duration is standard.

    \item \textbf{Overlapping TTA is widespread.}
    Top notebooks use 2.5s-stride sliding windows (yielding 23 windows per
    60s file) rather than non-overlapping 5s segments, effectively doubling
    the prediction density before aggregation.

    \item \textbf{Temporal smoothing is standard.}
    Gaussian smoothing ($\sigma$=1.0) or weighted-average kernels across
    adjacent segments reduce segment-boundary noise. Our v6 ablation
    confirmed $+0.006$ LB gain from smoothing.

    \item \textbf{Temperature scaling is common} ($T$=1.15--1.5).
    Softening predictions before ensemble averaging or threshold application
    improves calibration for AUC. We have not yet explored this.

    \item \textbf{Foundation models (Perch v2) are competitive baselines.}
    Perch v2 provides strong zero-shot predictions for 203/234 species.
    The BirdTransform notebook fuses CNN predictions with Perch embeddings.

    \item \textbf{Diverse backbones appear in top solutions.}
    ConvNeXt-V2, EfficientNetV2, Transformers, and EfficientNet-B2 are all
    represented, consistent with the 2025 finding that backbone diversity
    is critical for ensemble robustness.

    \item \textbf{No public notebook uses iterative self-training.}
    This is likely because it requires substantial compute and careful tuning,
    but the 2025 winners demonstrated it is the single most impactful technique
    when properly implemented (with noise injection, LB validation, and
    multi-backbone diversity across iterations).
\end{enumerate}

\paragraph{Implications for future experiments.}
The gap between our 0.811 and competitive 0.93+ likely requires:
(a) aligning training to 5s chunks with overlap TTA,
(b) multi-backbone ensemble diversity,
(c) proper iterative self-training with noise injection (mixup ratio=1,
stochastic depth) and LB-validated iterations, and
(d) inference-time calibration (temperature scaling, temporal smoothing).
Each of these should be validated independently on LB before combining.

% ------------------------------------------------------------------
\subsection{Experiment 16: Soft AUC Loss}
\label{sec:exp16}

\paragraph{Hypothesis.}
Directly optimizing the AUC ranking metric via a differentiable pairwise
surrogate loss should outperform BCE, which optimizes per-sample calibration
rather than ranking. The 2025 3rd-place solution reported success with this approach.

\paragraph{Setup.}
Soft AUC loss replaces BCE: for each class $c$, positive predictions
$\hat{y}^+$ and negative predictions $\hat{y}^-$ are compared pairwise:
\[
\mathcal{L}_\text{soft-AUC} = \frac{1}{|P||N|} \sum_{i \in P} \sum_{j \in N}
(1 - \sigma(\hat{y}_i - \hat{y}_j))
\]
Implemented efficiently using positive/negative partition and broadcast.
Applied to exp14's pipeline (5s chunks, mel augmentation, B0 backbone).
Training: 15 epochs, lr=$10^{-3}$, cosine schedule, 2-fold CV.

\paragraph{Results.}
\begin{table}[h]
\centering
\caption{Experiment 16 results (Soft AUC loss).}
\label{tab:exp16}
\begin{tabular}{lcccc}
\hline
\textbf{Fold} & \textbf{Best Epoch} & \textbf{val\_auc} & \textbf{ss\_auc} & \textbf{$\Delta$ vs exp14} \\
\hline
0 & 13 & 0.9831 & 0.7495 & $-0.033$ \\
1 & 9 & 0.9768 & 0.7204 & $-0.048$ \\
\hline
\end{tabular}
\end{table}

\paragraph{Analysis.}
Soft AUC loss \textbf{degraded} performance by 0.03--0.05 ss\_auc relative
to BCE (exp14). The ss\_auc trajectory showed an oscillating pattern, peaking
early and declining, suggesting training instability.

The failure likely stems from the multi-label sparse setting: with 234 classes
and typically 1--2 positives per sample, the pairwise loss is dominated by
easy negative pairs (233 negatives vs 1 positive), providing weak gradient
signal for discriminating between similar species. BCE, despite not directly
optimizing AUC, provides stronger per-class gradients that are more effective
for sparse multi-label tasks.

\textbf{Finding:} Soft AUC loss is not effective for sparse multi-label
classification with many classes. BCE with focal weighting (exp8) remains
the best loss function for this task.

% ------------------------------------------------------------------
\subsection{Experiment 17: ConvNeXt-V2 + GEM Pooling}
\label{sec:exp17}

\paragraph{Hypothesis.}
A top public notebook achieves high LB scores using ConvNeXt-V2 Tiny
(28M params) with 3-channel mel-spectrogram input and Generalized Mean (GEM)
frequency pooling. Replicating this architecture with our training recipe
and 234-class label mapping should yield a strong base model.

\paragraph{Key architecture differences from exp14.}
\begin{enumerate}
    \item Backbone: ConvNeXt-V2 Tiny (28M params) vs EfficientNet-B0 (5.3M)
    \item Input: 3-channel (mono expanded to RGB) vs 1-channel
    \item Normalization: per-sample mean/std vs min-max
    \item Pooling: GEM over frequency vs mean pool
    \item Head: Conv1d attention SED vs Linear attention SED
    \item Mel computation: torchaudio on GPU vs librosa precomputed
\end{enumerate}

\paragraph{Results.}
Six training attempts were made, systematically isolating failure modes:

\begin{table}[h]
\centering
\caption{Experiment 17: ConvNeXt-V2 training attempts.}
\label{tab:exp17}
\small
\begin{tabular}{llcc}
\hline
\textbf{Attempt} & \textbf{Configuration} & \textbf{val\_auc (E5)} & \textbf{ss\_auc} \\
\hline
1 & ImageNet pretrained, lr=$5 \times 10^{-4}$ & 0.500 & 0.48 \\
2 & ImageNet pretrained, lr=$5 \times 10^{-5}$ & 0.500 & 0.49 \\
3 & Public competition weights, full load & 0.500 & 0.50 \\
4 & Public weights + double-pow GEM fix & 0.500 & 0.50 \\
5 & Public weights, backbone only (heads random) & 0.500 & 0.48 \\
6 & ImageNet pretrained, standard GEM, per-epoch sched & 0.500 & 0.52 \\
\hline
\end{tabular}
\end{table}

All attempts produced val\_auc $\approx$ 0.50 (random chance), despite training
loss dropping rapidly to $\sim$0.027, indicating the model learns to predict
all-zeros (optimal for sparse multi-label BCE when class imbalance is extreme).

\paragraph{Root cause analysis.}
The failure has multiple contributing factors:

\textbf{(1) Extreme class imbalance without compensating sampler.}
With 234 classes and typically 1 positive per sample, the model minimizes BCE
by predicting near-zero for all classes. Our exp14 pipeline uses a
WeightedRandomSampler that oversamples rare species; exp17's unweighted
DataLoader provides no such balance. With ConvNeXt-V2's 28M parameters (5.3$\times$
B0), the model has more capacity to learn the trivial ``predict all zeros'' solution.

\textbf{(2) Species index mismatch with public weights.}
The public training notebook maps 206 species from \texttt{train.csv} to indices
0--205, while we map 234 species from \texttt{taxonomy.csv} to indices 0--233.
These are different alphabetical orderings due to the 28 extra species, making
the classification heads of the public weights meaningless in our framework.

\textbf{(3) A bug in the public training notebook.}
The public notebook sets \texttt{CFG.num\_classes = 206} as a class attribute,
but creates the model with \texttt{CFG()} which uses the dataclass default of 234.
The model thus has 234-class heads but is trained with 206-class label indices,
leaving 28 dimensions untrained.

\textbf{Finding:} ConvNeXt-V2 requires architecture-specific training adaptations
(class-balanced sampling, per-sample normalization, proper weight initialization)
that are not trivially portable from EfficientNet pipelines. The larger model
capacity makes it more susceptible to the all-zeros collapse mode.

% ------------------------------------------------------------------
\subsection{Experiment 20: Inference Post-Processing Evaluation}
\label{sec:exp20}

\paragraph{Hypothesis.}
The inference post-processing techniques from top public notebooks---overlapping
windows, time-shift TTA, temporal smoothing, and global soundscape prior---can
substantially improve predictions from our exp14 base model without retraining.

\paragraph{Setup.}
Using exp14 2-fold weights (ss\_auc=0.78), we evaluate each post-processing
technique incrementally on the full soundscape evaluation set (739 segments,
75 species with AUC-evaluable classes):

\begin{enumerate}
    \item \textbf{Overlapping windows}: 2.5s stride (50\% overlap), yielding
    23 windows per 60s file, merged back to 12 standard 5s segments by averaging
    overlapping predictions
    \item \textbf{Time-shift TTA}: shift each window by 1.25s and average
    original + shifted predictions
    \item \textbf{Temporal smoothing}: confidence-sharpened kernel
    $[0.05, 0.15, 0.60, 0.15, 0.05]$ with $p=1.5$ power sharpening
    \item \textbf{Global soundscape prior}: add 5\% of file-maximum probability
    to all segments (species detected anywhere in the file has elevated baseline)
\end{enumerate}

\paragraph{Results.}
\begin{table}[h]
\centering
\caption{Experiment 20: incremental post-processing evaluation (ss\_auc).}
\label{tab:exp20}
\begin{tabular}{lcc}
\hline
\textbf{Configuration} & \textbf{ss\_auc} & \textbf{$\Delta$ vs baseline} \\
\hline
exp14 baseline (no post-processing) & 0.7951 & --- \\
+ TTA only & 0.7996 & $+0.005$ \\
+ Overlap windows + TTA & 0.8136 & $+0.019$ \\
+ Overlap + TTA + Temporal smoothing & 0.8252 & $+0.030$ \\
+ Overlap + TTA + Gaussian $\sigma$=1.0 & 0.8238 & $+0.029$ \\
\textbf{+ Overlap + TTA + Smoothing + Global prior} & \textbf{0.8291} & $+\mathbf{0.034}$ \\
\hline
exp14+exp15 ensemble + full pipeline & 0.8021 & $+0.007$ \\
\hline
\end{tabular}
\end{table}

\paragraph{Analysis.}
The post-processing pipeline provides a cumulative $+0.034$ ss\_auc gain,
decomposed as follows:

\textbf{(1) Overlapping windows provide the largest single gain ($+0.014$).}
The 2.5s stride ensures that every bird vocalization is captured fully within
at least one window, eliminating segment-boundary truncation effects. This
is particularly important for short calls ($<$2.5s) that may straddle
the boundary between adjacent 5s segments.

\textbf{(2) Temporal smoothing adds $+0.011$.}
The confidence-sharpened kernel propagates high-confidence detections to
adjacent segments while suppressing noise. The sharpening power ($p=1.5$)
ensures that only confident predictions ($>$0.5) contribute meaningfully
to neighbors, preventing smoothing from diluting weak signals.

\textbf{(3) Global soundscape prior adds $+0.004$.}
If a species is detected with high confidence anywhere in the 60s recording,
its baseline probability is elevated by 5\% across all segments. This captures
the ecological prior that species presence is correlated across time within
a single recording.

\textbf{(4) Ensemble with exp15 hurts.}
Adding exp15 (self-training) weights to the ensemble \textbf{degraded}
performance from 0.8291 to 0.8021. This confirms our earlier finding (exp10)
that self-trained models overfit to train\_soundscapes distribution and
hurt generalization, even in ensemble with the base model.

\textbf{Finding:} Inference post-processing yields $+0.034$ ss\_auc at zero
training cost. The full pipeline (exp14 + overlap + TTA + smoothing + global
prior) achieves ss\_auc = 0.8291, our best local result. This should be
submitted to verify LB improvement over the 0.811 baseline.

% ------------------------------------------------------------------
\subsection{Experiment 17: ConvNeXt-V2 + GEM Pooling (Negative Result)}
\label{sec:exp17}

\paragraph{Motivation.}
The highest-scoring public training notebook uses ConvNeXt-V2 Tiny
(fcmae\_ft\_in22k\_in1k\_384) with 3-channel mel input and Generalized Mean (GEM)
frequency pooling, achieving LB scores above 0.90 in combination with
post-processing. We attempted to replicate this architecture.

\paragraph{Setup.}
ConvNeXt-V2 Tiny backbone (28.6M params), 3-channel mel spectrogram (mono expanded
to RGB), per-sample standardization ($\mu/\sigma$ normalization), GEM frequency
pooling with learnable exponent $p$ (init=3.0), attention SED head.
Training: 10 epochs, AdamW (lr=$5\times10^{-5}$), cosine schedule, batch size 32,
SpecAugment, mixup ($\alpha=0.4$), WeightedRandomSampler for class balance.

\paragraph{Results.}
Despite 8 training runs with progressive debugging, all runs collapsed to
val\_auc $\approx$ 0.50 (random). Root cause analysis revealed that
\texttt{model.forward\_features()} in timm's ConvNeXt-V2 \textbf{skips the final
LayerNorm2d} in the classifier head. The un-normalized features (max difference 0.627
from normalized output) caused the downstream GEM pooling and attention head to
receive unstable inputs, leading to all-zeros collapse. Even after correcting to
\texttt{model()} (which includes LayerNorm), the model still collapsed to val\_auc=0.50,
suggesting deeper architectural or hyperparameter incompatibility.

\begin{table}[h]
\centering
\caption{Experiment 17 ConvNeXt-V2 training runs (all failed).}
\begin{tabular}{lllr}
\toprule
Run & Change & val\_auc & ss\_auc \\
\midrule
1--4 & Initial attempts, public weights & 0.50 & -- \\
5 & Backbone-only weight loading, differential LR & 0.50 & -- \\
6 & ImageNet pretrained, standard GEM, lr=5e-5 & 0.50 & -- \\
7 & + WeightedRandomSampler & 0.50 & 0.47 \\
8 & + LayerNorm fix (forward\_features$\rightarrow$model) & 0.50 & 0.50 \\
\bottomrule
\end{tabular}
\end{table}

\textbf{Finding:} ConvNeXt-V2 replication failed after 8 attempts. The root cause
involves both the LayerNorm omission in \texttt{forward\_features()} and likely
additional factors (learning rate sensitivity, training data volume, or missing
self-training that the public notebook implicitly relies on). This negative result
motivated a strategic pivot to the Perch v2 foundation model approach.

% ------------------------------------------------------------------
\subsection{Strategic Pivot: Google Perch v2 Foundation Model}
\label{sec:perch_pivot}

Analysis of the BirdCLEF+ 2026 leaderboard (top score 0.952) and public notebooks
revealed that the dominant approach uses Google's \textbf{Perch v2} pre-trained bird
audio embedding model rather than training mel-spectrogram CNNs from scratch.

\paragraph{Key insight.}
Perch v2 is a foundation model trained on 14,795 bird species. It provides
1536-dimensional embeddings and species logits from 5-second audio windows.
Rather than learning audio representations from scratch (our exp3--17 approach),
the top solutions use Perch as a frozen feature extractor and train lightweight
classifiers (LogisticRegression, MLP, ProtoSSM) on its embeddings.

\paragraph{Public notebook analysis.}
The highest-voted public notebook family (\texttt{pantanal-distill}, 578 votes)
uses the pipeline:
\texttt{Perch} $\rightarrow$ \texttt{ProtoSSM + MLP} $\rightarrow$ \texttt{ensemble}
$\rightarrow$ \texttt{ResidualSSM} $\rightarrow$ \texttt{TTA + post-processing}.
A simpler Perch baseline (LB 0.910) uses:
\texttt{Perch} $\rightarrow$ \texttt{Bayesian priors} $\rightarrow$
\texttt{PCA + LogReg probes} $\rightarrow$ \texttt{Gaussian smoothing}.

We adopted the simpler baseline first for rapid LB validation, with plans to
upgrade to the full ProtoSSM pipeline.

% ==================================================================
\subsection{Experiments 21--25: Diagnosing the LB 0.910 Pipeline}
\label{sec:exp21-25}

The Perch~v2 baseline reached public LB 0.910 (exp20 v1), but we lacked any
local proxy aligned with that score. Submission v2 (PCA=64, $C{=}0.5$,
$\alpha{=}0.45$, 3$\times$ TTA) timed out, exhausting four of the five
remaining submission slots without yielding new information. Rather than
spend the remaining slots on blind variants, we ran a five-experiment
diagnostic sprint (exp21--25) on a single-machine CPU cluster to isolate
which components of the v1 pipeline contribute robust generalization and
which are fragile memorization.

A single Perch cache (708 windows, 1536-d embedding + 234-class score per
window, computed on the 59 fully-labeled train soundscapes) was shared
across all five experiments to amortize the 60~min Perch pass.
All experiments use the same 5-fold GroupKFold-by-site split (8 unique
sites: S03, S08, S13, S15, S18, S19, S22, S23) for OOF evaluation. Where
applicable, in-sample (no holdout) numbers are reported as a proxy for the
LB-style memorization upside.

\subsubsection{Experiment 21: OOF Ablation Harness}
\label{sec:exp21}

\paragraph{Setup.}
We re-implemented the v1 pipeline as a sequence of cumulative stages
(A through F) and evaluated each under both \emph{OOF} (5-fold GroupKFold
by site, prior tables fit on $N{-}1$ sites per fold) and \emph{in-sample}
(no holdout, prior tables fit on all 59 files) protocols.

\begin{table}[h]
\centering
\caption{Experiment 21: OOF vs in-sample macro-AUC for the v1 pipeline.}
\label{tab:exp21}
\begin{tabular}{llcc}
\hline
\textbf{Stage} & \textbf{Component added} & \textbf{OOF} & \textbf{In-sample} \\
\hline
A & Raw Perch v2 logits                              & 0.7390 & 0.7390 \\
B & + Genus proxy for Amphibia (max of same-genus)   & 0.7390 & 0.7390 \\
C & + Bayesian site$\times$hour priors (texture/event blend) & \textbf{0.488} & \textbf{0.977} \\
D & + Texture-taxa temporal smoothing ($\alpha{=}0.35$) & 0.485 & 0.978 \\
E & + PCA(32) LogReg probes ($C{=}0.25$, balanced)   & 0.539 & 0.979 \\
F & + Gaussian temporal kernel ($\sigma{=}1.0$)      & 0.524 & 0.978 \\
\hline
\end{tabular}
\end{table}

\paragraph{Analysis.}
Stage~C is the inflection point. Adding Bayesian priors raises in-sample
macro-AUC from 0.739 to 0.977 ($+0.238$) but collapses OOF from 0.739 to
0.488 ($-0.251$). The two protocols agree on raw Perch (stage~A), then
diverge by 0.49~AUC after prior fusion. Per-taxa breakdown (Table~\ref{tab:exp21-taxa})
shows the OOF damage concentrated in Amphibia and Insecta sonotypes
(\texttt{47158sonXX}, alphanumeric Insecta IDs).

\begin{table}[h]
\centering
\caption{Experiment 21 per-taxa OOF macro-AUC (fold-averaged).}
\label{tab:exp21-taxa}
\begin{tabular}{lccc}
\hline
\textbf{Taxa} & \textbf{A: raw} & \textbf{C: + priors} & \textbf{E: + probes} \\
\hline
Aves (mapped birds)            & 0.890 & 0.803 & 0.824 \\
Mammalia                       & 0.945 & 0.875 & 0.906 \\
Amphibia (mixed mapping)       & 0.804 & 0.524 & 0.599 \\
Insecta sonotypes (47158sonXX) & 0.500 & 0.108 & 0.162 \\
Reptilia                       & 0.500 & 0.341 & 0.743 \\
\hline
\end{tabular}
\end{table}

For the unmapped Insecta sonotypes (no Perch logit, no genus proxy), raw
Perch is uninformative (AUC=0.500). The Bayesian prior is then the
\emph{only} signal; under OOF site holdout the (site, hour) cells specific
to the held-out sites are absent, so only the global and hour-only
fall-backs fire. At unseen sites these fall-backs are anti-correlated with
truth, dragging AUC below chance to 0.108. The prior fusion is in effect
a site-conditioned look-up table that scores 0.977 in-sample because every
test row's site has been seen at training time.

\paragraph{Implication.}
LB 0.910 is consistent with the in-sample row of stage~F (0.978), not the
OOF row (0.524). The submission's leaderboard score reflects the test
soundscapes overlapping heavily with sites already seen during prior
fitting. If the hidden private test set contains genuinely novel sites,
the public-LB-optimized v1 pipeline is a candidate for severe shake-up.
Every subsequent experiment is therefore reported with both OOF
(generalization) and in-sample (memorization upside) numbers.

\subsubsection{Experiment 22: train\_audio Probe Transfer}
\label{sec:exp22}

\paragraph{Hypothesis.}
\texttt{train\_audio} contains 35{,}549 cleanly-recorded clips covering
234 species with reliable hard labels. If we extract Perch~v2 embeddings
on these clips and train a one-vs-rest linear classifier, the resulting
234-class probe should transfer to soundscape windows and either replace
or complement the per-class probes that exp21 trains on the 59-file
labeled subset.

\paragraph{Setup.}
First 5\,s of each \texttt{train\_audio} file passed through Perch~v2 to
obtain a $35{,}549 \times 1536$ embedding matrix (cached). Per-class
\texttt{LogisticRegression} ($C{=}0.1$, \texttt{class\_weight=balanced},
\texttt{liblinear}) trained one-vs-rest with the file's
\texttt{primary\_label} as positive. Of 234 classes, 192 had at least one
positive sample; 42 (mostly Insecta sonotypes) have no \texttt{train\_audio}
representation. Standardization fitted on \texttt{train\_audio} only.

Evaluation on the exp21 cache (708 SS windows, 71 active classes): the
$\texttt{train\_audio}$ probe alone, raw Perch alone, and uniform blends
$w \in \{0, 0.2, \ldots, 1.0\}$ where the final score is
$w \cdot \text{probe} + (1{-}w) \cdot \text{raw\_Perch}$.

\paragraph{Results.}
\begin{table}[h]
\centering
\caption{Experiment 22: train\_audio probe transfer to soundscape windows.}
\label{tab:exp22}
\begin{tabular}{lc}
\hline
\textbf{Configuration} & \textbf{SS macro-AUC} \\
\hline
Raw Perch alone ($w{=}0$)                    & 0.7390 \\
\texttt{train\_audio} probe alone ($w{=}1$)  & \textbf{0.6775} \\
Blend $w{=}0.2$ (best uniform)               & 0.7407 \\
Blend $w{=}0.4$                              & 0.7365 \\
Blend $w{=}0.5$                              & 0.7320 \\
Blend $w{=}0.8$                              & 0.7105 \\
Smart fill (probe used only for Perch-silent classes) & 0.7390 \\
\hline
\end{tabular}
\end{table}

\paragraph{What ``train\_audio fails'' means.}
The probe trained on 35{,}549 cleanly-recorded clips is, by itself,
\emph{worse} than the zero-shot Perch logit on the same evaluation
(0.677 vs 0.739). Adding it to Perch helps only at the smallest tested
weight ($w{=}0.2$ yields $+0.0017$ macro-AUC), and any larger contribution
degrades the ensemble. The ``smart blend'' variant---where the probe is
used \emph{only} for the 28 classes whose Perch logit is identically
zero (unmapped, no genus proxy)---produced exactly 0.7390, i.e.\ no
measurable change. The probe contributes nothing for the precise classes
it was supposed to rescue.

The mechanism is a clean$\to$field domain shift inside Perch's
embedding space:
\begin{itemize}
\item \texttt{train\_audio} clips are short ($\sim$5--30\,s),
single-species, foreground-dominant, often Xeno-Canto recordings made by
recordists who position the microphone close to a calling individual.
SNR is high and background acoustic energy is low.
\item Soundscape windows are 5\,s slices from passive 1-minute recorders
deployed in the Pantanal. Multiple species vocalize concurrently against
constant biophony (insects, frogs, distant species) and abiotic noise
(rain, wind, geophony).
\end{itemize}
A linear probe fitted to the first regime learns a hyperplane that
separates well-isolated species in the high-SNR Perch embedding; that
hyperplane sits in a region of embedding space the field windows rarely
visit. The result: the linear scores assigned to soundscape windows are
poorly ranked, and adding them to Perch dilutes Perch's already
domain-aware logits.

The 28-class ``smart fill'' result is the more diagnostic finding. For
the unmapped Insecta sonotypes, Perch returns a literal zero, so any
non-trivial score the probe assigns must be informative for AUC ranking.
That the smart-fill macro-AUC equals raw-Perch macro-AUC to four decimal
places means the probe's predictions on those classes are uncorrelated
with the truth---chance-level on average---confirming that the
\texttt{train\_audio} clips provide \emph{no transferable signal} for
the very classes that most need rescue.

\paragraph{Verdict.}
\texttt{train\_audio} probes in their current form are not a usable
lever. To make them work would require closing the clean$\to$field gap
\emph{before} probe training: e.g.\ Noisy-Student-style augmentation
(mixup with soundscape backgrounds, ESC-50 ambient noise, time-frequency
masking) on \texttt{train\_audio} embeddings; or a domain-adaptation
finetune of Perch's last layer with adversarial site-confusion. We
defer this to exp31+.

\subsubsection{Experiment 23: Pseudo-Labeling on Partial Soundscapes}
\label{sec:exp23}

\paragraph{Hypothesis.}
The 206 train soundscape files split into three populations: 59 files
with all 12 windows labeled, some files with a subset of windows
labeled, and the remainder with no labels at all. Pseudo-labeling the
\emph{partially-labeled} files would expand the probe training set
without introducing new acoustic distributions.

\paragraph{Reality.}
Of the 206 train soundscapes, the file-level label-coverage histogram is
in fact bimodal: 59 files at full coverage and \textbf{7 files} with
intermediate coverage; the remaining 140 files have zero windows
labeled. The pseudo-labeling pool is therefore $7 \times 12 = 84$
windows---two orders of magnitude smaller than the 17{,}892 pool used in
exp7's mel-CNN pipeline.

\paragraph{Setup.}
Perch extraction on the 7 partial files (84 windows; 8\,s wall time,
cached). Teacher predictions are raw Perch logits. Pseudo-labels:
top-$K{=}3$ classes per window with logit threshold 0; any real labels
present in the partial annotations are unioned in. Probe training:
\texttt{LogisticRegression} on PCA(32) of standardized embedding,
sample weight 1.0 on real (Y\_FULL) rows and 0.3 on pseudo rows.
Final score: $0.5 \cdot \text{raw Perch} + 0.5 \cdot \text{probe}$.

\paragraph{Results.}
\begin{table}[h]
\centering
\caption{Experiment 23: pseudo-labeling effect on OOF macro-AUC.}
\label{tab:exp23}
\begin{tabular}{lcc}
\hline
\textbf{Configuration} & \textbf{OOF AUC} & \textbf{$\Delta$ vs raw} \\
\hline
Raw Perch (baseline)                                  & 0.7390 & --- \\
+ Real-only probes (no priors, no pseudo)             & \textbf{0.8158} & $+0.0768$ \\
+ Real + pseudo probes (84 pseudo windows, 283 labels) & 0.8234 & $+0.0844$ \\
\hline
\end{tabular}
\end{table}

\paragraph{The accidental finding.}
The pseudo-label gain over real-only is small ($+0.0076$): with only 84
pseudo windows out of 708 evaluation rows there is little additional
signal. The unexpected result is the \emph{real-only probes} row:
adding PCA(32) probes to raw Perch \emph{without any prior fusion}
lifts OOF macro-AUC by $+0.077$ to 0.816. Compare this against
exp21's stage~E (probes \emph{after} prior fusion), which scored 0.539:
the same probe machinery produces a net $-0.20$ when fed prior-fused
input and a net $+0.077$ when fed raw input.

The probes were always doing useful work; the priors were destroying
the signal they had to work with. This experiment changes the
interpretation of stages~C--F in Table~\ref{tab:exp21}: priors are
not just a fragile addition on top of a solid base, they actively
delete the embedding-resident generalization that probes can otherwise
exploit.

\paragraph{Verdict.}
Pseudo-labeling on the available 7 partial files yields a small
positive contribution that we will keep for free, but the headline
result is the priors-removal recipe: \emph{raw Perch + PCA(32) LogReg
probes (no Bayesian priors)} hits OOF 0.816, our best site-honest
number to date.

\subsubsection{Experiment 24: Texture-Taxa Dedicated Head}
\label{sec:exp24}

\paragraph{Hypothesis.}
Insecta and Amphibia (``texture taxa'') are sustained calls with strong
file-level structure; the v1 PCA(32)+scalar probe features may be too
coarse for them. A richer feature set---the full 1536-d Perch
embedding, the file-mean embedding (averaging all 12 windows of the
same file as context), and the file-max-pooled Perch logits over the
mapped texture classes---should give the probe more to work with.

\paragraph{Setup.}
63 texture classes (those with \texttt{class\_name} $\in$
\{Amphibia, Insecta\}); 42 are active in the evaluation set. Per-class
\texttt{LogisticRegression} fitted in two configurations:
\begin{itemize}
\item \textbf{v1 probe} (control): PCA(32) of standardized embedding plus
the raw Perch logit for the target class. $C{=}0.25$, balanced.
\item \textbf{Rich probe}: full 1536-d standardized embedding $+$ 1536-d
standardized file-mean embedding $+$ standardized file-max-pooled
Perch logits over mapped texture classes (concatenated). $C{=}0.05$
(stronger regularization to compensate for higher dimensionality),
balanced.
\end{itemize}
Final score is $0.5 \cdot \text{raw Perch} + 0.5 \cdot \text{probe}$ in
both cases. 5-fold GroupKFold-by-site OOF.

\paragraph{Results.}
\begin{table}[h]
\centering
\caption{Experiment 24: texture-taxa-only macro-AUC (42 active classes).}
\label{tab:exp24}
\begin{tabular}{lc}
\hline
\textbf{Configuration} & \textbf{Texture macro-AUC} \\
\hline
Raw Perch                                  & 0.6399 \\
v1 probe (PCA(32)+scalar)                  & \textbf{0.7849} ($+0.145$) \\
Rich probe (full emb + file-mean + context)& 0.7226 ($+0.083$) \\
\hline
\end{tabular}
\end{table}

\begin{table}[h]
\centering
\caption{Experiment 24 by class\_name (mean of active per-class AUC).}
\label{tab:exp24-taxa}
\begin{tabular}{lccc}
\hline
\textbf{Class} & \textbf{Raw Perch} & \textbf{v1 probe} & \textbf{Rich probe} \\
\hline
Amphibia & 0.846 & 0.843 & 0.825 \\
Insecta  & 0.500 & 0.745 & 0.653 \\
\hline
\end{tabular}
\end{table}

\paragraph{Analysis.}
The richer feature set \emph{loses} to the minimal v1 probe by
$-0.062$ macro-AUC. Both taxa contribute to the gap: Amphibia drops
$-0.018$, Insecta drops $-0.092$. With a median of $\sim$10 positive
windows per active texture class across 5 folds, the per-class linear
classifier sees as few as 1--2 positives per training fold. A 1536+1536
dimensional feature vector with 1--2 positives is severely under-determined;
the strong $C{=}0.05$ regularization is not enough to prevent overfitting
to which-fold-the-positives-fell-into. The PCA(32) features already
contain the dominant directions of variation; further dimensions add
noise faster than signal at this sample size.

A second factor, visible only in the per-class breakdown, is that the
Insecta sonotypes that the v1 probe partially rescues
(\texttt{47158son07}, \texttt{47158son17}) are exactly the classes where
file-mean context is least informative: they are sustained tonal calls
whose 12 windows of a single file are nearly identical, so the file-mean
adds essentially the same vector again. The rich probe's extra features
are redundant rather than complementary.

\paragraph{Verdict.}
Keep the v1 PCA(32)+scalar probe head for texture taxa. Do not over-feature
small-N classes. Future texture-specific work should focus on supplying
\emph{more positive examples} (cross-year reuse of 2025 Pantanal data,
expert-curated sonotype clips) rather than richer features per example.

\subsubsection{Experiment 25: Domain-Shift / Site-Bias Diagnostics}
\label{sec:exp25}

\paragraph{Hypothesis.}
The dramatic OOF--in-sample gap in exp21 implies Perch embeddings
contain site-discriminative information. We measure how strong this
signal is and whether it justifies a domain-adaptation fine-tune
(iVAE / DANN / CORAL / per-site whitening) in exp30+.

\paragraph{Diagnostics.}
Five tests on the exp21 cache (708 windows, 8 sites, 13 unique hours):
\begin{description}
\item[D1] \textbf{Adversarial site classifier.} Multinomial LogReg
($\ell_2$, \texttt{lbfgs}) on the standardized 1536-d embedding
predicting the site label, evaluated by 5-fold StratifiedKFold accuracy.
Compared against chance ($1/8 = 0.125$) and against an analogous
hour-of-day classifier.
\item[D2] \textbf{Per-site leave-one-out OOF.} For each site $s$, hold
out all windows belonging to $s$ and report raw-Perch macro-AUC on the
held-out subset.
\item[D3] \textbf{Texture-taxa site concentration.} For each texture
class, count how many of the 8 sites contribute at least one positive
window, and compute a Gini-like dispersion $G = 1 - \sum_s p_s^2$ where
$p_s$ is the fraction of positives at site $s$. $G{=}0$ means
single-site, $G \to 1$ means uniform.
\item[D4] \textbf{Embedding silhouette by site.} Standardize, project to
PCA(20), compute \texttt{silhouette\_score} treating site as the cluster
label.
\item[D5] \textbf{Per-class fold-AUC variance.} 5-fold GroupKFold-by-site
raw-Perch per-class AUC, report mean and std across folds for the 71
active classes.
\end{description}

\paragraph{Results.}

\begin{table}[h]
\centering
\caption{Experiment 25: domain-shift diagnostic summary.}
\label{tab:exp25}
\begin{tabular}{lr}
\hline
\textbf{Diagnostic} & \textbf{Value} \\
\hline
D1: site classifier accuracy           & 0.999 \\
\quad chance ($1/8$)                   & 0.125 \\
\quad ratio above chance               & $\boldsymbol{8.0\times}$ \\
D1: hour classifier accuracy           & 0.992 (chance 0.077, $13\times$) \\
D2: per-site raw Perch macro-AUC mean  & 0.666 \\
\quad std / min / max                  & 0.141 / 0.487 (S18) / 0.829 (S22) \\
D4: silhouette\_score(site, PCA(20))   & $+0.096$ \\
D5: cross-fold per-class AUC std (mean / max) & 0.012 / 0.184 \\
\hline
\end{tabular}
\end{table}

\begin{table}[h]
\centering
\caption{Experiment 25 D3: texture classes with single-site positives.}
\label{tab:exp25-d3}
\begin{tabular}{lcrr}
\hline
\textbf{Class} & \textbf{Taxa} & \textbf{N pos} & \textbf{N sites} \\
\hline
\texttt{47158son02} & Insecta  & 7  & 1 \\
\texttt{47158son06} & Insecta  & 18 & 1 \\
\texttt{47158son14} & Insecta  & 12 & 1 \\
\texttt{47158son17} & Insecta  & 43 & 1 \\
\texttt{47158son19} & Insecta  & 5  & 1 \\
\texttt{47158son21} & Insecta  & 22 & 1 \\
\texttt{25073}      & Amphibia & 12 & 1 \\
\texttt{25092}      & Amphibia & 24 & 1 \\
\texttt{326272}     & Amphibia & 23 & 1 \\
\texttt{517063}     & Amphibia & 307& 1 \\
\hline
\multicolumn{4}{l}{\emph{(15+ texture classes total exhibit single-site positives.)}} \\
\end{tabular}
\end{table}

\paragraph{Analysis.}

\textbf{D1: the embedding is strongly site-confounded.} A linear classifier
on Perch embeddings recovers the recording site with near-perfect
accuracy (0.999, 8$\times$ chance). The embedding is encoding---through
its 1536 dimensions---enough of the recorder's frequency response,
ambient biophony, and microclimate to make sites linearly separable.
The hour-of-day signal is also recoverable (0.992 vs $0.077$ chance,
$13\times$), partly because each site was visited at consistent hours
in this dataset (a measurement-protocol confound, not a Perch
limitation). A randomized stratified holdout cannot detect either
because every site/hour is in both train and validation; only a
group-aware split (such as our GroupKFold-by-site) exposes the
problem.

\textbf{D2: per-site difficulty is highly non-uniform.} Raw-Perch macro-AUC
ranges from 0.487 (S18) to 0.829 (S22), $\sigma{=}0.141$ across 6
evaluable sites (S15 and S23 are NaN because no class is active in
their hold-out subset). S18 sits below chance, meaning that on that
site Perch's class ranking is anti-correlated with the truth---a sign
that the species mixture and ambient profile of S18 differ enough
from Perch's training distribution that its top logits become
systematically wrong. A pipeline that performs well only on the
``easy'' sites can post a misleadingly high LB if the public test
sample is dominated by them.

\textbf{D3: texture sonotypes are single-site labels.} The Insecta
sonotype IDs (\texttt{47158sonXX}) and several Amphibia codes have all
of their positive windows concentrated at one site. This is the
binding constraint that no domain-adaptation method can lift: there
is literally no cross-site supervision to learn from. Removing the
site signal from the embedding (CORAL / DANN) or modelling a site-free
latent (iVAE) would not help, because for these classes the only
training-time decision boundary that can possibly separate them from
negatives also separates the site they were recorded at from the
others. Their AUC$=$0.500 in the most-stable cell of D5 is mathematical,
not algorithmic.

\textbf{D4: silhouette is small-positive.} A silhouette of $+0.096$
(within the $[-1, 1]$ range, where 0 means no cluster structure) is
not a tight clustering, but it is clearly positive: the embedding
arranges windows into loose site-shaped neighborhoods in PCA(20)
space. Combined with D1's near-perfect linear separability, this
suggests the site signal lies along high-rank directions that linear
methods can exploit but PCA's first 20 components only partly capture.

\textbf{D5: cross-fold variance is moderate on average, very high for
some Amphibia classes.} The mean across-fold std of per-class AUC is
0.012 (low), but the worst classes have std up to 0.184. The top of
the high-variance list is dominated by Amphibia
(\texttt{22961}, \texttt{24321}, \texttt{1491113}, \texttt{517063},
\texttt{22967}); the most-stable list is dominated by classes whose
AUC is stuck at exactly 0.500 across all folds (the single-site
texture sonotypes from D3, plus a handful of Reptilia). ``Stable''
in this context means uniformly uninformative.

\paragraph{Implication for exp30+.}
Domain adaptation is justified for the 71 active classes that have
multi-site labels: D1 shows the site signal is large enough to act on,
and D2 shows the OOF cost is large enough to motivate the work. Linear
candidates---per-site mean subtraction, CORAL alignment, simple
adversarial debiasing of the probe stage---are cheap and should be
tried first. Heavier methods (iVAE, DANN finetune of Perch's last
layer) require GPU and a larger labeled corpus; they should only
follow if the linear methods plateau.

For the 15+ single-site texture sonotypes, no embedding-side method
can manufacture cross-site supervision. The realistic options are:
(a) augment with 2025 Pantanal soundscape clips containing the same
sonotype IDs, (b) collect expert-curated reference clips per sonotype
and treat them as one-shot prototypes (a ProtoNet-style head over
Perch), or (c) accept AUC$=$0.500 on these classes and rely on the
remaining 56 active classes to drive macro-AUC.

\subsubsection{Synthesis: Robust Recipe for exp30 v1}
\label{sec:exp21-25-synthesis}

The five experiments converge on a single robust recipe that we will
submit as exp30 v1.

\begin{table}[h]
\centering
\caption{exp21--25 cumulative effect on OOF macro-AUC (708 SS windows, 71 active classes).}
\label{tab:exp21-25-summary}
\begin{tabular}{llcc}
\hline
\textbf{Recipe} & \textbf{Source} & \textbf{OOF} & \textbf{In-sample} \\
\hline
Raw Perch v2                         & exp21 A & 0.7390 & 0.7390 \\
v1 (priors + probes + smoothing)     & exp21 F & 0.524  & 0.978 \\
Probes only (no priors)              & exp23   & \textbf{0.8158} & --- \\
Probes only + 84-window pseudo       & exp23   & \textbf{0.8234} & --- \\
\hline
\end{tabular}
\end{table}

\textbf{Drop} the Bayesian prior fusion stage entirely. Under Val-B
it costs $-0.20$ vs raw Perch and is naively interpretable as site
memorization. exp27 reveals a more nuanced reason: priors do carry
real signal under the test-realistic Val-A regime ($+0.06$ over raw),
but the PCA-projected embedding fed into the probe already encodes
site identity, so explicit prior fusion double-counts the same signal
and adds noise. The PP design therefore loses to R1 (probes only) in
\emph{both} regimes (Val-A $-0.108$, Val-B $-0.289$).

\textbf{Keep} PCA(32) + per-class \texttt{LogisticRegression}
($C{=}0.25$, balanced) probes blended 50/50 with raw Perch
(exp23 establishes this is where the $+0.077$ OOF gain comes from).

\textbf{Do not} replace the probe head with rich 1536-d features for
texture taxa: the v1 minimal features beat the rich variant by
$-0.062$ macro-AUC because per-class positive counts are too small
(exp24).

\textbf{Add} the 84-window pseudo-label pool from the 7 partial
soundscape files for an extra $+0.0076$ OOF (exp23). It is free.

\textbf{Do not} blend in a \texttt{train\_audio} probe: the
clean$\to$field domain shift makes the probe worse than zero-shot
Perch and adds nothing for the 28 Perch-silent classes (exp22).

\textbf{Plan} domain adaptation for exp31+: site signal is real
($8\times$ chance, exp25 D1) and per-site OOF varies by 0.34 AUC
(exp25 D2), but linear methods (per-site whitening, CORAL) should
be tried before heavy ones (DANN / iVAE finetune).

\textbf{Predicted exp30 v1 OOF (Val-B)}: $\approx 0.82$. \textbf{Predicted
LB}: unknown until exp27 introduces a second evaluation regime that
matches the actual test conditions. The exp21--25 OOF protocol
(GroupKFold by site) is the worst-case Val-B regime where every
validation site is unseen at training time. exp27 shows the public
test set is closer to a Val-A regime (seen sites, new files), under
which the recipe ranking changes substantially.

\subsection{Experiment 26: Robust Recipe Variant Sweep}
\label{sec:exp26}

\paragraph{Motivation.} exp23 discovered that dropping prior fusion
lifts Val-B macro-AUC from 0.524 to 0.816. exp26 sweeps nine recipes
over the same cached features to (a) confirm the lever, (b) probe
adjacent variants (per-file embedding centering, weak priors at
varying shrinkage $k$, Gaussian temporal smoothing), and (c) pick a
canonical recipe for the next submission slot.

\paragraph{Recipes.} All built on the cached exp21 Perch features
(708 windows $\times$ 234 classes). Each recipe is evaluated under
5-fold GroupKFold by site (Val-B):
\begin{description}
\item[R0] Raw Perch v2.
\item[R1] PCA(32) + LogReg probe ($C{=}0.25$, balanced),
50/50 blend with raw Perch. Reproduces exp23.
\item[R2a--d] R1 with weak priors fused into the base, with
shrinkage $k_h{=}k_s{\in}\{50, 100, 200, 500\}$
($k_{sh}$ scaled accordingly).
\item[R3] R1 + per-file embedding centering (subtract the
file-mean embedding from each window before PCA).
\item[R4] R3 + best-$k$ priors.
\item[R5] R1 + 5-tap Gaussian temporal smoothing across the
12 windows of each file ($\sigma{\approx}1.0$).
\item[R6] R3 + best-$k$ priors + Gaussian smoothing.
\end{description}

\begin{table}[h]
\centering
\caption{Experiment 26: Val-B macro-AUC for nine recipe variants
(708 windows, 71 active classes).}
\label{tab:exp26}
\begin{tabular}{llc}
\hline
\textbf{Recipe} & \textbf{Description} & \textbf{Val-B AUC} \\
\hline
R0  & Raw Perch v2                         & 0.7390 \\
R1  & Probes only (no priors)              & 0.8148 \\
R2a & Probes + weak priors $k{=}50$        & 0.5137 \\
R2b & Probes + weak priors $k{=}100$       & 0.5095 \\
R2c & Probes + weak priors $k{=}200$       & 0.5086 \\
R2d & Probes + weak priors $k{=}500$       & 0.5066 \\
R3  & R1 + per-file centering              & 0.7486 \\
R4  & R3 + best-$k$ priors                 & 0.4633 \\
\textbf{R5} & \textbf{R1 + Gaussian smoothing} & \textbf{0.8190} \\
R6  & R3 + best-$k$ priors + Gaussian      & 0.4472 \\
\hline
\end{tabular}
\end{table}

\paragraph{Findings.} \textbf{R5 wins.} Adding a 5-tap Gaussian smoother
across the 12 windows of each file lifts R1 from 0.8148 to 0.8190
($+0.0042$) at zero compute cost. \textbf{Weak priors do not survive
GroupKFold-by-site at any shrinkage.} Even at $k{=}500$ (heavy
shrinkage toward the global rate), Val-B drops to 0.507 because
validation sites are entirely absent from the prior tables and the
hour-only fallback is uninformative across sites. \textbf{Per-file
centering (R3) is a wash}: it lifts Reptilia from 0.500 to 0.856 by
removing file-level DC offset, but kills Insecta from 0.693 to 0.519
because single-site sonotype embeddings ARE the file-mean for those
classes. The Insecta loss dominates the macro average.

\paragraph{Per-taxa breakdown.} Comparing R0, R1, R3, and R6 by taxa
(Table~\ref{tab:exp26-taxa}) shows the structural risk of stacked
priors: R6 collapses every taxon, with Insecta dropping below 0.10
because the per-file centering destroys the only signal that worked
for those classes. R1 is the only recipe that lifts every active
taxon over R0.

\begin{table}[h]
\centering
\caption{Experiment 26: per-taxa mean Val-B AUC.}
\label{tab:exp26-taxa}
\begin{tabular}{lcccc}
\hline
\textbf{Taxon} & \textbf{R0} & \textbf{R1} & \textbf{R3} & \textbf{R6} \\
\hline
Aves      & 0.890 & 0.905 & 0.898 & 0.752 \\
Mammalia  & 0.945 & 0.946 & 0.948 & 0.814 \\
Amphibia  & 0.846 & 0.841 & 0.825 & 0.491 \\
Insecta   & 0.500 & 0.693 & 0.519 & 0.075 \\
Reptilia  & 0.500 & 0.762 & 0.856 & 0.282 \\
\hline
\end{tabular}
\end{table}

\paragraph{Verdict.} R5 is the canonical Val-B recipe. We carry it
forward to exp27 to test whether it also wins under a more
test-realistic evaluation regime.

\subsection{Experiment 27: Dual Validation (Val-A + Val-B)}
\label{sec:exp27}

\paragraph{Motivation.} The 2025-style GroupKFold-by-site protocol
(Val-B) holds out entire sites. Public discussion in the BirdCLEF
2026 forum (lb928 thread, Tucker Arrants) points out that this is
\textit{overly harsh} for this competition because the actual test
soundscapes overlap training sites. A model that holds out entire
sites at validation time cannot learn the legitimate site-level
species priors that exist at test time. The recommendation is to
evaluate under \textbf{both} regimes:
\begin{description}
\item[Val-A] file-stratified-by-site 5-fold (sites seen at train
time, files held out). Matches the actual test regime.
\item[Val-B] GroupKFold by site (entire site held out). Hedges
against the unseen-site fraction in private LB.
\end{description}

\paragraph{Method.} For each recipe we compute Val-A AUC (round-robin
file assignment within each site, weighted toward Val-A by $\geq{1}$
per site) and Val-B AUC (5 disjoint site partitions, exp21--26
protocol). Five recipes are evaluated:
\begin{description}
\item[R0] Raw Perch v2.
\item[PR] Raw + Bayesian prior fusion (no probes).
\item[R1] Probes only, 50/50 blend with raw Perch.
\item[PP] Priors + probes (the LB 0.910 design).
\item[R5] R1 + Gaussian temporal smoothing.
\end{description}

\begin{table}[h]
\centering
\caption{Experiment 27: Val-A (seen-site, new-file) vs Val-B
(unseen-site) macro-AUC. The $\Delta$ column is Val-A $-$ Val-B.}
\label{tab:exp27}
\begin{tabular}{lccc}
\hline
\textbf{Recipe} & \textbf{Val-A} & \textbf{Val-B} & \textbf{$\Delta$ A$-$B} \\
\hline
R0 raw Perch                & 0.7390 & 0.7390 & $\phantom{-}0.0000$ \\
PR priors only              & 0.8006 & 0.5121 & $+0.2884$ \\
\textbf{R1 probes only}     & \textbf{0.8886} & \textbf{0.8372} & $+0.0514$ \\
PP priors $+$ probes        & 0.7805 & 0.5482 & $+0.2324$ \\
\textbf{R5 R1 $+$ Gaussian} & \textbf{0.8922} & \textbf{0.8418} & $+0.0504$ \\
\hline
\end{tabular}
\end{table}

\paragraph{Mechanism: why probes already contain priors.} Before
interpreting the table, it helps to be precise about what each
component is doing. The Bayesian \textbf{prior} is a lookup table
indexed by site and hour. For each (site, hour) bin in the training
data, we compute the empirical fraction of windows containing each
species, convert to a logit, and add $\lambda$ times that logit to
Perch's raw score. Two consequences: (a) the prior consults only
metadata (filename-derived site and hour); it never examines the
audio of the test window; (b) every test window in the same
(site, hour) bin receives the \emph{same} additive bias on every
class. The prior cannot distinguish between a quiet and a loud
window, or between two adjacent 5-second windows of the same file.

The \textbf{probe} is a per-class binary logistic regression trained
on PCA(32) features of Perch's 1536-d embedding. Crucially, the
embedding encodes acoustic site identity automatically: recordings
from S22 sound systematically different from those of S03 (microphone
characteristics, background fauna, ambient noise floor), and these
differences survive the PCA projection because they are dominant
variance directions. When LogReg learns ``does species $k$ appear?''
from training labels, it discovers any discoverable correlation in
the 32-d feature space, including the correlation
``S22-shaped embeddings tend to have species $k$.'' In other words,
the probe learns the same site$\to$species association the prior
table makes explicit, except (i) it learns it from data rather than
from a hand-coded lookup, and (ii) its predictions vary
\emph{per window} because each window has its own embedding, whereas
the prior's contribution is constant across all windows in a bin.

This explains every row of Table~\ref{tab:exp27}. PR (priors only)
beats raw Perch on Val-A ($+0.062$) because the lookup is accurate
when the test site is in the training table. R1 (probes only) beats
PR by another $+0.088$ because the probe captures the same site
signal \emph{plus} per-window acoustic discrimination. PP (priors
$+$ probes) loses to R1 in both regimes because the prior adds a
bin-constant bias on top of the probe's already-site-aware
prediction; this bias does not improve discrimination between windows
in the same bin and dilutes the probe's per-window signal in
proportion to $\lambda$. Val-B kills priors entirely because the
table lookup misses: held-out sites are absent from
\texttt{site\_to\_i}, and the hour-only fallback is uninformative
across sites. The probe degrades more gracefully on Val-B
(0.889 $\to$ 0.837) because it can still match acoustic patterns of
species that occur at multiple sites, even when no training
embeddings come from the held-out site.

\paragraph{Reframe of the prior story.} The exp21--26 narrative was:
``priors give an in-sample lift of $+0.20$ that vanishes in OOF; they
are a memorization shortcut.'' The exp27 result complicates that:
\begin{enumerate}
\item Priors \textbf{do} provide a real Val-A lift over raw Perch
($+0.062$), confirming Tucker's point that site-level species
priors carry signal when test sites overlap training sites.
\item However, probes provide a \textbf{larger} Val-A lift ($+0.150$)
over raw Perch, and an even larger Val-A advantage over
priors-only ($+0.088$). The PCA-projected embedding implicitly
encodes site identity, so the LogReg probe captures the same
``which species at which site'' association that the prior table
captures, plus per-window discriminative information.
\item Combining the two (PP) is \textbf{worse} than probes alone in
both regimes ($-0.108$ Val-A, $-0.289$ Val-B). The two pathways
double-count the site signal and contribute interference rather
than complementary information.
\end{enumerate}

\textbf{R5 wins both regimes} (Val-A 0.892, Val-B 0.842). It is the
canonical recipe.

\paragraph{Caveats.}
\begin{itemize}
\item Val-A folds 3 and 4 are S22-heavy because S22 owns 39 of 59
labeled files. Recipes that exploit S22-specific structure may
benefit unfairly. The recipe ranking is, however, consistent
across folds 0--2 (which are more balanced).
\item The PP simulation here uses PCA-only probe input, while the
LB 0.910 PP submission uses richer per-class features
(\texttt{raw\_col}, \texttt{prior\_col}, sequence statistics).
The qualitative ordering R1 $>$ PP is robust across this
modeling choice; the exact gap may differ.
\item Val-A still does not perfectly match the actual test regime.
The competition test set may include both seen and unseen sites
in unknown proportion. Val-A and Val-B together bracket the
likely range; a recipe that wins both is preferable to one that
maximizes either alone.
\end{itemize}

\paragraph{Implication for exp30.} The next submission slot,
\texttt{birdclef-2026-exp20-to-exp30} v1, runs the R1 recipe with
priors fully disabled (\texttt{LAMBDA\_*}=0, PCA-only probe input,
fixed 0.5 blend). Based on Table~\ref{tab:exp27}, the predicted
public LB is $0.88$--$0.93$, with the upside scenario (LB $>$ 0.910)
arising precisely because R1 captures site signal more efficiently
than PP. v2 will be R5 (R1 + Gaussian smoothing).

% ------------------------------------------------------------------
\subsection{Experiment 28--30: Beyond Perch-only (\S Kaggle top-scorer grounding)}
\label{sec:exp28_30}

\paragraph{Motivation.} Public Kaggle discussion (digested 2026-04-20)
confirms that top single-model scores (Boredom 0.947, Ali Ozan 0.925,
Salman 0.922) are achieved \emph{without} self-distillation or
pseudo-labeling on unlabeled soundscapes. The reported decomposition
(\citeauthor{hengck23}) is
\[
\text{LB}_{0.93} = \underbrace{\text{Perch2} + \text{proxy class} + \text{temporal} + \text{PCA}}_{\approx 0.92} + \text{prior} + \text{state space} + \text{per-class norm}.
\]
Our exp20 baseline (LB 0.910) already implements proxy class, PCA,
and temporal sequence features via
\texttt{seq\_features\_1d}. The remaining cheap levers are
\emph{per-class rank normalisation} and \emph{state space smoothing};
the harder but necessary lever for a top-family score is a
second-family model (SED CNN). Pseudo-labeling is deprioritised
(Don Mathis's public post reports failure; top five submissions skip
it).

\paragraph{exp28 --- Perch probe v2 (CPU).}
Rebaseline the starter-notebook-frozen probe (PCA 64, $C{=}0.50$,
$\alpha{=}0.40$) against our exp20 v1 config (PCA 32, $C{=}0.25$)
under Val-A/Val-B. Add per-class rank normalisation, and sweep
smoothers (Gauss $\sigma\in\{0.5,1,1.5\}$, 1-D Kalman, HMM with
transition counts from \texttt{train\_soundscapes\_labels.csv}).

\paragraph{exp29 --- HGNetV2-B0 SED (Salman's recipe).}
Backbone \texttt{hgnetv2\_b0.ssld\_stage2\_ft\_in1k},
20-second training clips, raw-waveform mixup, CE loss on the sum of
clipwise and framewise-max logits, 20 epochs of AdamW with cosine
schedule. Training data is \texttt{train\_audio} + labeled
soundscapes only. The recipe is a direct re-implementation of a
publicly disclosed top-20 solo recipe; an independent reproduction
serves as both a diversity model for the final ensemble and a
validity check on the claim.

\paragraph{exp30 --- ConvNeXt-small SED and ensemble.}
Same recipe with a transformer-style backbone for architectural
diversity (EliKal reports that Perch-KD makes backbones
interchangeable; we test the converse claim in a non-KD setting).
The final submission is the Val-A-optimal blend of Perch probe v2
$+$ HGNetV2-B0 SED $+$ ConvNeXt-small SED, subject to a hard
constraint that Val-B does not regress by more than $0.01$.
ONNX export of the SED models verifies that the end-to-end
pipeline fits inside the Kaggle 90-minute CPU budget.

Results are deferred to Section~\ref{sec:exp28_30_results}; the
experiments run on an RTX 5090 with a fixed 12-hour wall budget
starting 2026-04-20.

\subsection{Experiment 34--36: Blend ceiling and SED sweep (negative results)}
\label{sec:exp34_36}

\paragraph{exp34 --- Perch~$\oplus$~SED29 z-score blend.}
Z-score normalised per-class fusion of Perch probe v2 scores
(Val-A 0.8943) and exp29 HGNetV2-B0 SED scores (Val-A 0.7374) at
mixing weight $\alpha{=}0.80$. A 1-D Gaussian smoother of
$\sigma{=}0.5$ is applied to the blended logits. The blend reaches
\textbf{Val-A 0.9062, Val-B 0.8503}. Platt and isotonic calibration
per fold both \emph{reduce} Val-A, because the per-fold positive count
is too small (typically $<$10) to fit a reliable link function.

\paragraph{exp35 --- ConvNeXtV2-tiny SED (diversity probe).}
Identical Salman recipe on a ConvNeXtV2-tiny backbone. Best
val\_auc $0.6468$ at epoch~17 with decline at epochs~18--20, confirming
convergence. Pearson correlation between exp29 and exp35 on active
classes is $0.231$ --- low enough to indicate architectural diversity
--- yet the 3-way simplex search assigns weight $w_{35}{=}0$: the
absolute AUC is too low to survive the simplex, so diversity alone
is not sufficient.

\paragraph{exp36 --- Regularised SED sweep.}
Three configs sharing the exp29 recipe with additional regularisation:
SpecAugment ($F{=}24$, $T{=}80$, $N{=}2$ each, $P{=}0.7$), label
smoothing $0.1$, weight decay $5{\times}10^{-2}$, and early-stopping
patience of $5$ at up to $25$ epochs. Backbones are HGNetV2-B0 (A),
HGNetV2-B1 (B), and EfficientNet-B0
(\texttt{tf\_efficientnet\_b0.ns\_jft\_in1k}, C). Single-model Val-A
scores are $0.628/0.607/0.647$, all \emph{below} the unregularised
exp29 baseline of $0.737$. The 3-way simplex of
Perch$+$SED29$+$SED$_\text{new}$ again collapses to weight $0$ on the
new model in all three cases. We conclude that stacking SpecAugment,
label smoothing, and aggressive weight decay simultaneously
over-regularises the SED head; future re-tries must introduce
regularisers one at a time.

\paragraph{Leaderboard: v6 with SED blend and priors on.}
The exp34 recipe is deployed in the \texttt{exp20-to-exp30} submission
notebook with the v1 0.910 priors-on configuration restored (prior
fusion is required to avoid random predictions on unmapped classes in
the hidden test set). With post-blend Gaussian smoothing $\sigma{=}0.5$
and batch\_files$=16$ (validated to fit in the Kaggle 90-minute CPU
budget), the submission achieves \textbf{public LB 0.912} on
2026-04-20, an improvement of $+0.002$ over the v1 Perch-only
baseline. The local gain (Val-A $+0.012$) transfers only partially to
the leaderboard, suggesting that the prior fusion already covers much
of the diversity that SED29 contributes on unmapped classes.

\subsection{Experiment 37: Pivot to public ProtoSSM pipeline (LB 0.927)}
\label{sec:exp37_protossm}

\paragraph{Motivation.}
After the SED sweep ceiling of LB 0.912 was established, we reviewed
publicly shared top-scoring notebooks. The hengck23 decomposition
(\texttt{example\_and\_discussion/lb928.md}) identifies the missing
components as \emph{state-space smoothing} and \emph{per-class
normalisation}, both of which we had tested either as post-hoc Kalman
(negative, exp30) or rank-norm (negative, exp28). We pivoted to
evaluate whether a \emph{trained} state-space head on Perch embeddings
--- as distinct from post-hoc smoothing --- would yield a larger gain.

\paragraph{Pipeline.}
The reference implementation is the community-shared ``ProtoSSM v5''
notebook, forked as
\texttt{notebooks/birdclef-2026-perch-distill}. The pipeline is:
Perch v2 embeddings $\to$ selective state-space model (Mamba-style,
$d_\mathrm{model}{=}320$, $n_\mathrm{layers}{=}4$,
$d_\mathrm{state}{=}32$) with prototypical classification head and
cross-attention $\to$ per-class MLP probes
(hidden sizes $[256,128]$, PCA-128 features) $\to$ fusion with
OOF-optimised mixing weight ($0.50$ in submit mode) $\to$ Residual
SSM (second-pass boosting on first-pass errors, $d{=}128$, 2 layers,
correction weight $0.35$) $\to$ per-taxon temperature (Aves $1.10$,
texture $0.95$) $\to$ file-level top-$k$ scaling ($k{=}2$) $\to$
rank-aware scaling (power $0.4$) $\to$ delta-shift temporal smoothing
($\alpha{=}0.2$) $\to$ per-class threshold sharpening. All heads are
trained in-notebook at submit time on the labeled soundscapes;
inference on hidden test runs Perch once and performs the subsequent
heads on top of the 1536-d embeddings.

\paragraph{Runtime engineering.}
Two earlier versions (v3, v4) timed out on the 90-minute CPU budget.
Diagnosis from the kernel log localised the cost to two factors.
First, the cache of Perch embeddings on the labeled soundscapes
(\texttt{jaejohn/perch-meta} dataset) failed to mount at the
notebook-hardcoded path \texttt{/kaggle/input/perch-meta}; the
competition environment actually mounts datasets at
\texttt{/kaggle/input/datasets/<owner>/<slug>}, adding approximately
$560$~s (9 minutes) of Perch re-computation on the $107$ labeled
soundscape files. Second, \texttt{batch\_files$=16$} --- a value we
briefly tried for memory-pressure reasons --- increased per-file
Perch cost compared to \texttt{batch\_files$=32$} due to poorer
XLA compilation amortisation over the larger hidden test.

\paragraph{Fix.}
We extended \texttt{resolve\_full\_cache\_paths} to enumerate four
candidate mount paths
(\texttt{/kaggle/input/datasets/jaejohn/perch-meta},
\texttt{/kaggle/input/perch-meta},
\texttt{/kaggle/input/jaejohn/perch-meta}, and the user-configured
\texttt{full\_cache\_input\_dir}), and reverted \texttt{batch\_files}
to $32$. The updated dev run (v6) reports
\texttt{Loading cached full-file Perch outputs from\dots} at wall time
$113.7$~s, completes in $280.4$~s for a 20-file dry run with 2 of 2
auxiliary models trained, and produces a well-formed submission of
shape $(240, 235)$.

\paragraph{Result.}
\textbf{Public LB 0.927}, an improvement of $+0.015$ over the exp34
SED-blend submission (LB 0.912) and $+0.017$ over the v1 Perch-only
baseline (LB 0.910). This result validates the hengck23 decomposition:
the missing component on our pipeline was not a second audio backbone
but a trained temporal model on the Perch embeddings. The per-class
normalisation and trained SSM together recover most of the reported
$0.924$--$0.930$ range for Perch-based solutions.

\paragraph{Lessons for our pipeline.}
The LB 0.927 result is obtained with essentially unmodified public
infrastructure; our pipeline-specific contributions (exp21--36) do
not appear in this submission. Remaining opportunities are
(i)~blending the exp29 SED stream into the ProtoSSM fusion step,
which our local analysis indicates should yield an additional
$0.002$--$0.010$ under appropriate z-score normalisation;
(ii)~back-porting the ProtoSSM and post-processing stack to
\texttt{exp20-to-exp30} so that our SED-blend asset remains active;
(iii)~systematic re-evaluation of per-class thresholds under the
new fusion, which the public notebook currently leaves at the
default $0.5$ in submit mode.

\subsection{Experiment 38: ONNX Perch and the cache-consistency failure}

\paragraph{Motivation.}
Public notebook
\texttt{birdclef-26-two-pass-ssm-advanced-pp} reports LB~0.933 using
ONNX Runtime for the Perch forward pass, quoting an approximate
$150\times$ speedup over the TensorFlow SavedModel. We asked
whether adopting ONNX in our ProtoSSM v5 pipeline frees enough
runtime budget to admit additional inference streams (e.g.\ the
SED29 blend that plateaued at LB~0.912 in exp34), without changing
the scoring pipeline.

\paragraph{Submissions (2026-04-21).}
Three controlled submissions isolated the effect of the runtime
swap on the LB while holding every other component of the v6
pipeline fixed.

\begin{center}
\begin{tabular}{lllr}
\toprule
Version & Training cache & Test inference & Public LB \\
\midrule
v6  & TF (\texttt{jaejohn/perch-meta}) & TF SavedModel & \textbf{0.927} \\
v9  & TF (\texttt{jaejohn/perch-meta}) & ONNX Runtime  & 0.924 \\
v10 & ONNX (recomputed in-notebook)    & ONNX Runtime  & 0.919 \\
\bottomrule
\end{tabular}
\end{center}

v9 replaces the test-time backbone with \texttt{perch\_v2.onnx} while
the training caches used to fit the MLP probes, ProtoSSM v5, and
ResidualSSM remained the TF-generated artefacts in
\texttt{jaejohn/perch-meta}. v10 additionally drops the TF cache and
recomputes every training-time Perch output with ONNX, matching the
end-to-end ONNX pipeline used by the public 0.933 notebook. The dev
wall time decreased from $280$~s (v6) to $170$~s (v9) to $358$~s (v10;
the additional time is spent re-generating the labelled-soundscape
cache from scratch at submit time).

\paragraph{Diagnosis.}
Our prior hypothesis --- that TF and ONNX Perch produce slightly
different embeddings and that a TF-fit probe head degrades when fed
ONNX features at test time --- predicts the ordering
$\text{v10} > \text{v9}$ once the train/test backbone is made
consistent. The LB results invert this prediction:
v10 (consistent ONNX) loses an additional $0.005$ relative to v9,
and $0.008$ relative to v6. Two mechanisms are compatible with this
ordering:
\begin{enumerate}
\item The external \texttt{jaejohn/perch-meta} cache is not a raw
Perch dump; it likely embeds pre-aggregation choices (proxy
reduction, windowing) that our in-notebook recomputation does not
exactly reproduce, so switching to an in-notebook ONNX cache both
changes the backbone and silently changes the training feature
distribution.
\item The \texttt{perch\_v2.onnx} export carries a small accuracy
cost relative to the TF SavedModel on our specific pipeline, and
the public LB~0.933 notebook only hits that figure because its
isotonic-calibration and adaptive-delta post-processing stack
masks the backbone loss.
\end{enumerate}
We cannot yet distinguish (1) from (2) from the LB alone.

\paragraph{Implication.}
ONNX is orthogonal to score under our current pipeline: its value
is runtime headroom, not accuracy. To realise the LB~0.93+ regime we
still need a second evidence stream (SED29 blend) and the
post-processing upgrades from the public notebook. The cache
recomputation path is a non-trivial source of noise and should be
avoided unless we adopt the public pipeline's post-processing at
the same time.

\subsection{Experiment 39: SED29 blend on the ProtoSSM pipeline}

\paragraph{Setup.}
Reverting to the v6 configuration (TF SavedModel, TF-generated
\texttt{jaejohn/perch-meta} cache, \texttt{batch\_files=32}) as a
known LB~0.927 baseline, we inject the exp29 SED29 stream into the
\texttt{perch-distill} post-processing pipeline. The blend is
identical to the one measured locally in exp34: a per-class
$z$-score blend with $\alpha_{\text{Perch}} = 0.80$ applied on the
logit-space fusion output, restricted to mapped active classes,
followed by Gaussian temporal smoothing with
$\sigma = 0.5$ across the 12-window axis. The SED model
(HGNetV2-B0 with attention head, trained in exp29) is loaded from
the existing \texttt{ultimatumgame/birdclef2026-model-weights}
dataset, producing 20-second windowed predictions that are
replicated across each 5-second Perch output row within a 20-second
block. No additional training is performed at submit time.

\paragraph{Expectation.}
Local measurement in exp34 gave Val-A $+0.012$ relative to the
raw Perch output at the same $\alpha$. Under the ProtoSSM fusion
the absolute level is higher ($0.927$ rather than $\sim 0.894$), and
the fusion already absorbs part of the SED signal through the
trained temporal head. We therefore expect a smaller transfer,
tentatively $+0.002$ to $+0.005$ LB, putting the target in the
$0.929$--$0.932$ range. The submission is scored as v11 in the
same kernel.

\subsection{Experiment 39: SED blend result on hybrid pipeline (LB 0.929)}

\paragraph{v11 failure mode.}
The first attempt at bundling SED29 into the LB~0.927 v6 configuration
(v11: TF SavedModel end-to-end plus SED) reached the submission
\texttt{COMPLETE} state with an empty \texttt{publicScore}, matching
the \emph{scoring failed} pattern observed in v8. Dev wall was $341$~s
for a 20-file dry run; extrapolated to the hidden-test regime the
Perch-TF forward dominates at roughly $95$~min for 600 files, which
exceeds the 90-minute CPU budget once ProtoSSM and ResidualSSM
training are accounted for. The root cause was runtime, not blend
correctness.

\paragraph{v12: hybrid pipeline.}
v12 keeps the TF-generated \texttt{jaejohn/perch-meta} cache so that
the MLP probes and ProtoSSM v5 are trained on the same feature
distribution that v6 saw, but routes hidden-test audio through
ONNX Perch. This combines the LB~0.927 training regime (jaejohn
cache) with the runtime headroom of ONNX (dev wall drops from
$341$~s to $238$~s, roughly $100$~s saved per 20 files, scaling to
a $\sim 50$~min saving over 600 files). Onto this we add the SED29
z-score blend ($\alpha_{\text{Perch}} = 0.80$, Gaussian smoothing
$\sigma = 0.5$) exactly as measured in exp34.

\paragraph{Result.}
\textbf{Public LB 0.929}, a new maximum. This is $+0.002$ over the
v6 TF-only baseline and $+0.005$ over the ONNX-only v9 run.

\paragraph{Decomposition.}
The four 2026-04-21 submissions isolate each component:

\begin{center}
\begin{tabular}{lllr}
\toprule
Version & Change & LB & $\Delta$ \\
\midrule
v6  & TF cache + TF test                         & $0.927$ & --- \\
v9  & TF cache + ONNX test                       & $0.924$ & $-0.003$ \\
v10 & ONNX cache + ONNX test                     & $0.919$ & $-0.008$ \\
v11 & TF cache + TF test + SED blend             & \emph{fail} & timeout \\
v12 & TF cache + ONNX test + SED blend (hybrid)  & $\mathbf{0.929}$ & $+0.002$ \\
\bottomrule
\end{tabular}
\end{center}

Reading by column:
\begin{itemize}
\item ONNX swap alone costs $-0.003$ (v6 vs v9). The public notebook's
$0.933$ figure therefore does not come from the ONNX backend itself;
it comes from the post-processing stack that the same notebook bundles.
\item Recomputing the training cache with ONNX (v10) deepens the
penalty to $-0.008$, confirming that the MLP/ProtoSSM heads are fit
to the exact feature distribution present in
\texttt{jaejohn/perch-meta} and that any in-notebook recomputation
introduces a silent distribution shift.
\item SED29 blend transfers $+0.005$ from dev to LB in v12, matching
the upper end of the exp34-based forecast of $+0.002$ to $+0.005$.
Local Val-A had reported $+0.012$ under the same $\alpha$; the
transfer ratio ($\approx 0.42$) is consistent with ProtoSSM already
absorbing part of the temporal signal that SED would otherwise
contribute.
\item The net $+0.002$ for v12 relative to v6 is the sum of the
SED gain ($+0.005$) and the ONNX penalty ($-0.003$). ONNX is a
pure enabler in this pipeline: without it, the SED blend cannot be
added because the runtime budget is already saturated (v11).
\end{itemize}

\paragraph{Interpretation.}
The hybrid recipe
--- TF-domain training features, ONNX test inference, SED blend
in logit space, post-blend Gaussian smoothing ---
is the configuration that yields the current LB maximum of $0.929$.
Every ablation we have tried so far drops the score, and the
specific combination of keeping the jaejohn cache while switching
the test backbone to ONNX is required to reproduce the result.

\paragraph{Remaining gap to $0.933$.}
The public \texttt{birdclef-26-two-pass-ssm-advanced-pp} pipeline,
which shares the Perch+ProtoSSM+ResidualSSM backbone with v12 but
reports LB~$0.933$, differs from ours in its post-processing: it
performs per-class isotonic calibration on OOF scores and
applies adaptive $\delta$-smoothing that blends uncertain windows
toward their neighbours (rather than the fixed-$\sigma$ Gaussian we
use). These two passes operate on the final fused logits after the
SED blend and are orthogonal to the decomposition above, so they
are the natural next component to port.
